\begin{center}
{\Large \bfseries \textcolor{lyred}{第十一章 曲线积分与曲面积分}}
\end{center}

\vspace{6pt}
{\large \bfseries \textcolor{lyred}{第11.1 节 对弧长的曲线积分}}

\vspace{4pt}
{\bfseries \textcolor{lyred}{一.第一类曲线积分的定义}}

\textcolor{lyred}{定义}.设L 为xOy 面内一段光滑曲线,
(
)
,
f x y 在L 上有界,将L 任意分成n 段
is
\Delta , 
弧长也记为
is
\Delta ,
(
)
,
i
i
is
\psi \eta 
\forall 
\in \Delta ,作和
(
)
,
n
i
i
i
i
f
s
\psi \eta 
=
\Delta 
\sum 
,若在无限细分L 的过程中, 
随
\{
\}
1max
i
i n
s
\lambda 
\le \le 
=
\Delta 
\to 
,该和总是趋向于同一个只依赖于
(
)
,
f x y 和L 的常数I ,则 
称I 为
(
)
,
f x y 在L 上\textcolor{lyred}{对弧长的曲线积分},记为
(
)
,
L
f x y ds
\int 
,若L 是闭曲线,也记为 
(
)
,
L
f x y ds
\int \nabla 
;称
(
)
,
f x y 为\textcolor{lyred}{被积函数}, L 为\textcolor{lyred}{积分弧段}. 
\textcolor{lyred}{物理意义}:设L 为曲线型材料,有连续密度
(
)
,x y
\mu 
,则
(
)
,
L
M
x y ds
\mu 
=\int 
. 
\textcolor{lyred}{几何意义}:设
(
)
,
f x y \ge 
,连续,则以L 为准线且平行z 轴的柱面介于曲面
z =
, 
(
)
,
z
f x y
=
之间部分的面积为
(
)
(
)
lim
,
,
n
i
i
i
i
L
A
f
s
f x y ds
\lambda 
\psi \eta 
\to 
=
=
\Delta 
=
\sum 
\int 
,例如: 
x
y
ax
a
x
y ds
+
=
-
-
\int \nabla 
为
x
y
ax
+
=
含在
x
y
z
a
+
+
=
内上半部分的面积. 
\textcolor{lyred}{定理}.设
(
)
,
f x y 在光滑曲线L 上连续,则
(
)
,
L
f x y ds
\int 
存在. 
\textcolor{lyred}{注}.若L 分段光滑,则规定
(
)
(
)
(
)
,
,
,
L
L
L
f x y ds
f x y ds
f x y ds
=
+
\int 
\int 
\int 
. 
\textcolor{lyred}{注}.对于分段光滑空间曲线\Gamma 和\Gamma 上的连续函数
(
)
, ,
f x y z ,也可以定义\textcolor{lyred}{对弧长的 }
\textcolor{lyred}{曲线积分}
(
)
(
)
, ,
lim
,
,
n
i
i
i
i
i
f x y z ds
f
s
\lambda 
\psi \eta 
\to 
=
\Gamma 
=
\Delta 
\sum 
\int 
. 
\vspace{4pt}
{\bfseries \textcolor{lyred}{二.第一类曲线积分的性质}}

\textcolor{lyred}{性质1}.
(
)
(
)
(
)
(
)
,
,
,
,
L
L
L
f x y
g x y
ds
f x y ds
g x y ds
\alpha 
\beta 
\alpha 
\beta 
\pm 
=
\pm 




\int 
\int 
\int 
; 
\textcolor{lyred}{性质2}.
(
)
(
)
(
)
,
,
,
L
L
L
L
f x y ds
f x y ds
f x y ds
+
=
+
\int 
\int 
\int 
; 
\textcolor{lyred}{性质3}. 1
L
ds
s
\cdots 
=
\int 
,一般地,
L
k ds
k s
\cdots 
=
\cdots 
\int 
; 
\textcolor{lyred}{性质4}.设在L 上
(
)
(
)
,
,
f x y
g x y
\le 
,则
(
)
(
)
,
,
L
L
f x y ds
g x y ds
\le 
\int 
\int 
; 
\textcolor{lyred}{推论}.
(
)
(
)
,
,
L
L
f x y ds
f x y ds
\le 
\int 
\int 
. 
\textcolor{lyred}{注}.也有类似于定积分的估值定理和积分中值定理. 
\textcolor{lyred}{例}.求
(
)
L
I
x
y
x
y ds
=
+
-
-
\int \nabla 
,其中
(
)
(
)
:
L
x
y
-
+
-
=
. 
解.
(
)
(
)
11 2
L
L
I
x
y
ds
ds
\pi 
\pi 


=
-
+
-
-
=
\cdots 
=
\cdots 
\cdots 
=


\int 
\int 
\nabla 
\nabla 
. 
\textcolor{lyred}{例}.求
I
y
z ds
\Gamma 
=
+
\int \nabla 
,其中
: x
y
z
y
x

+
+
=
\Gamma 
=

. 
解.
I
x
y
z ds
ds
\pi 
\Gamma 
\Gamma 
=
+
+
=
\cdots 
=
\int 
\int 
\nabla 
\nabla 
. 
\vspace{4pt}
{\bfseries \textcolor{lyred}{三.对称性}}

1.若积分曲线是平面曲线,则参照二重积分; 
2.若积分曲线是空间曲线,则参照三重积分. 
\textcolor{lyred}{例}.求
(
)
cos
x
L
I
x
y
ye
ds
=
+
+
\int \nabla 
,其中
:
L x
y
+
=. 
解.
4 2
L
I
ds
=
+
+
\cdots 
=
\int \nabla 
. 
\textcolor{lyred}{例}.求
L
x
y
I
ds
a
b


=
+




\int \nabla 
,其中
:
L x
y
R
+
=
. 
解.
(
)
L
L
x
y
I
ds
x
y
ds
R
a
b
a
b
a
b
\pi 






=
+
=
+
+
=
+












\int 
\int 
\nabla 
\nabla 
. 
\textcolor{lyred}{例}.设
(
)
:
x
L
y
+
-
=的周长为a ,求
(
)
L
I
x
y
ds
=
+
\int \nabla 
. 
解.
:
L x
y
y
+
=
,故
(
)
1 1
L
L
L
I
yds
y
ds
ds
a
=
=
-+
=
=
\int 
\int 
\int 
\nabla 
\nabla 
\nabla 
. 
\textcolor{lyred}{例}.求
(
)
L
I
y
x
x
y
x
y
ds


=
+
+
+
+


\int \nabla 
,其中
:
L x
y
x
+
=
. 
解.
(
)
(
)
(
)
2 2
L
L
L
I
x
x
x ds
xds
ds
\pi 


=
+
=
+
=
+
=
+


\int 
\int 
\int 
\nabla 
\nabla 
\nabla 
. 
\textcolor{lyred}{例}.求
(
)
I
xy
z ds
\Gamma 
=
+
\int \nabla 
,其中
:
x
y
z
x
y
z

+
+
=
\Gamma +
+
=

. 
解.
(
)
(
)
(
)
I
xyds
x
y ds
xy
yz
zx ds
x
y
z ds
\Gamma 
\Gamma 
\Gamma 
\Gamma 
=
-
+
=
+
+
-
+
+
=
\int 
\int 
\int 
\int 
\nabla 
\nabla 
\nabla 
\nabla 
(
)
(
)
x
y
z
x
y
z
ds
ds
\pi 
\Gamma 
\Gamma 


+
+
-
+
+
-
=-
=-


\int 
\int 
\nabla 
\nabla 
. 
\textcolor{lyred}{例}.求
(
)
I
x
y
z ds
\Gamma 
=
+
+
\int \nabla 
,其中
:
x
y
z
x
y

+
+
=
\Gamma +
=

. 
解.
(
)
I
xds
zds
x
y ds
\Gamma 
\Gamma 
\Gamma 
=
+
=
+
+
=
\int 
\int 
\int 
\nabla 
\nabla 
\nabla 
. 
\vspace{4pt}
{\bfseries \textcolor{lyred}{四.计算法}}

\textcolor{lyred}{定理}.设
()
()
: x
x t
L
y
y t
=

=

,
t
\alpha 
\beta 
\le \le 
,其中
()
x t
'
,
()
y t
'
均连续,并且
()
()
x
t
y
t
'
'
+
\neq 
, 
则
(
)
()
()
()
()
,
,
L
f x y ds
f
x t
y t
x
t
y
t dt
\beta 
\alpha 
'
'
=
+




\int 
\int 
. 
\textcolor{lyred}{推论}.(1)设
()(
)
:
L y
y x
a
x
b
=
\le 
\le 
,则
(
)
()
()
,
,
b
L
a
f x y ds
f
x y x
y
x dx
'
=
+




\int 
\int 
; 
特别地,
(
)
:
L y
y
a
x
b
=
\le 
\le 
平行于x 轴时,
(
)
(
)
,
,
b
L
a
f x y ds
f x y
dx
=
\int 
\int 
. 
(2)设
()(
)
:
L x
x y
c
y
d
=
\le 
\le 
,则
(
)
()
()
,
,
d
L
c
f x y ds
f
x y
y
x
y dy
'
=
+




\int 
\int 
; 
特别地,
(
)
:
L x
x
c
y
d
=
\le 
\le 
平行于y 轴时,
(
)
(
)
,
,
d
L
c
f x y ds
f x
y dy
=
\int 
\int 
. 
(3)设在极坐标下,
()
:
L \rho 
\rho \theta 
=
,\alpha 
\theta 
\beta 
\le 
\le 
,则 
(
)
()
()
()
()
,
cos ,
sin
L
f x y ds
f
d
\beta 
\alpha 
\rho \theta 
\theta \rho \theta 
\theta 
\rho 
\theta 
\rho 
\theta 
\theta 
'
=
+




\int 
\int 
. 
\textcolor{lyred}{例}.求
L
I
yds
=\int 
,其中L 为
y
x
=
上
(
)
0,0
O
与(
)
1,1
B
之间的一段弧. 
解.
(
)
(
)
5 5
I
x
x
dx
x
x dx
'
=
\cdots 
+
=
\cdots 
+
=
-
\int 
\int 
. 
\textcolor{lyred}{例}.求
(
)
x
y
L
I
x
y e
ds
+
=
+
\int \nabla 
,其中L 为
y
x
=
-
, y
x
=\pm 所围区域边界. 
解.
1 :
L
y
x
=-,
x
-
\le 
\le 
,
cos
:
sin
x
L
y
\theta 
\theta 
=


=

,
\pi 
\pi 
\theta 
\le 
\le 
,
3 :
L
y
x
=
,
x
\le 
\le 
, 
(
)
(
)
(
)
cos
sin
x
x
I
x
x e
dx
e d
x
x e
dx
e
\pi 
\pi 
\theta 
\theta 
\theta 
-
=
-
+
+
+
+
=
-
\int 
\int 
\int 
. 
\textcolor{lyred}{例}.求
L
I
y ds
=\int \nabla 
,其中
(
)
:
L
x
y
x
y
+
=
-
为双纽线. 
解.
:
cos2
L \rho 
\theta 
=
,故
sin
L
I
yds
d
\pi 
\rho 
\theta 
\rho 
\rho 
\theta 
'
=
=
\cdots 
+
=
\int 
\int 
(
)
(
)
(
)
4 sin
4 sin
4 sin
2 2
d
d
d
\pi 
\pi 
\pi 
\theta 
\rho 
\rho \rho 
\theta 
\theta 
\rho 
\rho 
\theta 
\theta \theta 

'
'
+
=
+
=
=
-




\int 
\int 
\int 
. 
\textcolor{lyred}{例}.求
x
y
ax
+
=
位于
x
y
z
a
+
+
=
内部分的面积. 
解.
L
A
a
x
y ds
=
-
-
\int \nabla 
,在极坐标下,
:
cos
L
a
\rho 
\theta 
=
,
\pi 
\pi 
\theta 
-
\le 
\le 
,故 
cos
A
a
d
a
a
ad
a
\pi 
\pi 
\pi 
\pi 
\rho 
\rho 
\rho 
\theta 
\theta 
\theta 
-
-
'
=
-
\cdots 
+
=
-
\cdots 
=
\int 
\int 
. 
\textcolor{lyred}{例}.求
x
y
+
=位于
x
y
z
+
+
=内部分的面积. 
解.
L
A
x
y ds
=
-
-
\int \nabla 
,其中
cos
:
sin
x
t
L
y
t
=

=

,0
t
\pi 
\le \le 
,故 
1 cos
sin
sin cos
3sin
cos
3sin cos
A
t
t
t
tdt
t
t
t
tdt
\pi 
\pi 
\pi 
=
-
-
\cdots 
=
\cdots 
=
\int 
\int 
. 
\textcolor{lyred}{注}.类似地,设光滑空间曲线
()
()
()
:
x
x t
y
y t
z
z t
=


\Gamma 
=

=

,
t
\alpha 
\beta 
\le \le 
,则 
(
)
()
()
()
()
()
()
, ,
,
,
f x y z ds
f
x t
y t
z t
x
t
y
t
z
t dt
\beta 
\alpha 
\Gamma 
'
'
'
=
\cdots 
+
+




\int 
\int 
. 
\textcolor{lyred}{例}.求
(
)
I
x
y
z ds
\Gamma 
=
+
+
\int 
,其中\Gamma 为(
)
1,1,1
A
到(
)
2,3,4
B
的直线段. 
解.
:
x
t
y
t
z
t
=+


\Gamma 
=
+

=
+

,
[
]
0,1
t \in 
,故
(
)
13 14
I
t
dt
=
+
+
+
=
\int 
. 
\textcolor{lyred}{例}.求
y
I
ds
x
y
z
\Gamma 
=
+
+
\int \nabla 
,其中
:
x
y
z
x
y
x

+
+
=
\Gamma 
+
=

. 
解.
y
I
ds
yds
\Gamma 
\Gamma 
=
=
\int 
\int 
\nabla 
\nabla 
,其中
(
)
1 cos
:
sin
2 1 cos
2sin 2
x
t
y
t
t
z
t

=+

\Gamma 
=


=
-
+
=

,0
t
\pi 
\le \le 
,故 
(
)
sin
1 cos
cos
cos
2 2
t
t
t
I
t
dt
d
\pi 
\pi 
=
\cdots 
+
=-
+
=
-
\int 
\int 
. 
\vspace{4pt}
{\bfseries \textcolor{lyred}{五.物理应用}}

设\Gamma 为曲线型材料,具有连续密度
(
)
, ,
x y z
\mu 
,则 
1.质心坐标
x ds
x
ds
\mu 
\mu 
\Gamma 
\Gamma 
=\int 
\int 
,
y ds
y
ds
\mu 
\mu 
\Gamma 
\Gamma 
=\int 
\int 
,
z ds
z
ds
\mu 
\mu 
\Gamma 
\Gamma 
=\int 
\int 
; 
2.对l 的转动惯量
lI
d
ds
\mu 
\Gamma 
=\int 
,其中(
)
, ,
d x y z 为(
)
, ,
x y z 到l 的距离; 
3.对(
)
,
,
x
y
z
处单位质点的引力 
(
)
(
)
(
)
,
,
x
x
y
y
z
z
F
G
ds
ds
ds
r
r
r
\mu 
\mu 
\mu 
\Gamma 
\Gamma 
\Gamma 


-
-
-
=




\int 
\int 
\int 
\rho 
. 
\textcolor{lyred}{例}.求半径为R ,中心角为2\alpha 的均匀圆弧关于对称轴的转动惯量. 
解.取x 轴为对称轴,原点为圆心建立坐标系,则
x
L
I
y
ds
\mu 
=
=
\int 
(
)
(
)
(
)
sin
sin
cos
sin
cos
R
R
R
d
R
\alpha 
\alpha 
\mu 
\theta 
\theta 
\theta 
\theta 
\mu \alpha 
\alpha 
\alpha 
-
-
+
=
-
\int 
. 
\textcolor{lyred}{补充练习 }
1.求
(
)
sin
L
I
x
x
y
x
y
y ds
=
+
+
+
-
\int \nabla 
,其中
(
)
:
L x
y
+
-
=. 
解.
(
)
L
L
L
I
x
y
y ds
yds
ds
\pi 
=
+
-
=
=
=
\int 
\int 
\int 
\nabla 
\nabla 
\nabla 
. 
2.求
L
I
x
y ds
=
+
\int \nabla 
,其中
:
L x
y
y
+
=-
. 
解.
:
L x
y
y
=\pm -
-
,
dx
y
dy
y
y
--
=\pm 
-
-
,
dy
ds
x dy
y
y
'
=
+
=
-
-
, 
故
2 2
4 2
dy
dy
I
y
y
y
y
y
-
-
-


=
-
\cdots 
=
=
+
=


+
-
-
\int 
\int 
; 
或者,
cos
:
1 sin
x
t
L
y
t
=


=-+

,故
2sin
sin
cos
t
t
I
t dt
dt
\pi 
\pi 
=
-
\cdots 
=
-
=
\int 
\int 
; 
或者,
:
2sin
L \rho 
\theta 
=-
,
ds
d
d
\rho 
\rho 
\theta 
\theta 
'
=
+
=
,故 
2sin
I
d
d
\pi 
\pi 
\rho 
\theta 
\theta 
\theta 
-
-
=
\cdots 
=
-
\cdots 
=
\int 
\int 
. 
\vspace{6pt}
{\large \bfseries \textcolor{lyred}{第11.2 节 对坐标的曲线积分}}

\vspace{4pt}
{\bfseries \textcolor{lyred}{一.变力沿曲线作功}}

设质点在变力
(
)
,
F
P Q
=
\rho 
的作用下沿有向曲线
\approx 
AB
\Gamma =
移动,将L 任意分成n 小段 
\dots 
i
i
i
L
M
M
-
=
,在每一小段上,变力作功
(
)
(
)
,
,
i
i
i
i
i
i
i
i
W
F
M
M
F
r
\psi \eta 
\psi \eta 
-
\approx 
\cdots 
=
\cdots \Delta 
\xi \xi \xi \xi \xi \xi \xi \rho 
\rho 
\rho 
\rho ,其中 
(
)
,
i
i
i
i
r
OM
x y
=
=
\xi \xi \xi \xi \rho 
\rho 
,故
(
)
(
)
(
)
lim
,
lim
,
,
n
n
i
i
i
i
i
i
i
i
i
i
i
W
F
r
P
x
Q
y
\lambda 
\lambda 
\psi \eta 
\psi \eta 
\psi \eta 
\to 
\to 
=
=
=
\cdots \Delta =
\cdots \Delta 
+
\cdots \Delta 




\sum 
\sum 
\rho 
\rho 
, 
其中
x
F 作的功为
(
)
lim
,
n
i
i
i
i
P
x
\lambda 
\psi \eta 
\to 
=
\cdots \Delta 
\sum 
,
y
F 作的功为
(
)
lim
,
n
i
i
i
i
Q
y
\lambda 
\psi \eta 
\to 
=
\cdots \Delta 
\sum 
,分别记为 
(
)
,
L
P x y dx
\int 
,
(
)
,
L
Q x y dy
\int 
,故
(
)
(
)
(
)
,
,
,
L
L
W
P x y dx
Q x y dy
F x y
dr
=
+
=
\cdots 
\int 
\int 
\rho 
\rho . 
\vspace{4pt}
{\bfseries \textcolor{lyred}{二.第二类曲线积分的定义}}

设
\approx 
L
AB
=
为光滑有向曲线,
(
)
,
f x y 在L 上有界,将L 任意分成n 段
\dots 
i
i
i
L
M
M
-
=
, 
其中
(
)
,
i
i
i
M
x y
=
,
(
)
,
i
i
iL
\psi \eta 
\forall 
\in 
,作和
(
)
,
n
i
i
i
i
f
x
\psi \eta 
=
\Delta 
\sum 
,若在无限细分L 的过程中, 
随着
\{
\}
1max
i
i n
s
\lambda 
\le \le 
=
\Delta 
\to 
,该和总是趋向于一个只依赖于
(
)
,
f x y 和L 的常数I ,则 
称I 为
(
)
,
f x y 在L 上\textcolor{lyred}{对坐标x 的曲线积分},记为
(
)
,
L
f x y dx
\int 
; 
称
(
)
,
f x y 为\textcolor{lyred}{被积函数}, L 为\textcolor{lyred}{积分弧段}. 
类似地,\textcolor{lyred}{对坐标y 的曲线积分}
(
)
(
)
,
lim
,
n
i
i
i
i
L
f
x y dy
f
y
\lambda 
\psi \eta 
\to 
=
=
\Delta 
\sum 
\int 
. 
设
(
)
,
F
P Q
=
\rho 
为
\approx 
L
AB
=
上的有界向量值函数,则它在L 上\textcolor{lyred}{对坐标的曲线积分}为 
L
L
F dr
Pdx
Qdy
\cdots 
=
+
\int 
\int 
\rho 
\rho 
,其中
(
)
,
dr
dx dy
=
\rho 
称为\textcolor{lyred}{有向曲线元素}. 
\textcolor{lyred}{定理}.设
(
)
,
f x y 在光滑有向曲线L 上连续,则它在L 上对坐标的曲线积分存在. 
\textcolor{lyred}{注}.设\Gamma 为空间光滑有向曲线,
(
)
, ,
f x y z 在\Gamma 上的有界函数,则 
(
)
(
)
, ,
lim
,
,
n
i
i
i
i
i
f x y z dx
f
x
\lambda 
\psi \eta 
\to 
=
\Gamma 
=
\Delta 
\sum 
\int 
,
(
)
(
)
, ,
lim
,
,
n
i
i
i
i
i
f x y z dy
f
y
\lambda 
\psi \eta 
\to 
=
\Gamma 
=
\Delta 
\sum 
\int 
, 
(
)
(
)
, ,
lim
,
,
n
i
i
i
i
i
f x y z dz
f
z
\lambda 
\psi \eta 
\to 
=
\Gamma 
=
\Delta 
\sum 
\int 
; 
设
(
)
,
,
F
P Q R
=
\rho 
,则F dr
Pdx
Qdy
Rdz
\Gamma 
\Gamma 
\cdots 
=
+
+
\int 
\int 
\rho 
\rho 
,其中
(
)
,
,
dr
dx dy dz
=
\rho 
. 
\textcolor{lyred}{物理意义}.质点在变力
(
)
,
,
F
P Q R
=
\rho 
作用下沿有向曲线
\approx 
AB
\Gamma =
移动,则所做的功 
W
F dr
Pdx
Qdy
Rdz
\Gamma 
\Gamma 
=
\cdots 
=
+
+
\int 
\int 
\rho 
\rho 
. 
\vspace{4pt}
{\bfseries \textcolor{lyred}{三.第二类曲线积分的性质}}

\textcolor{lyred}{性质1}. (
)
L
L
L
F
F
dr
F dr
F
dr
\alpha 
\beta 
\alpha 
\beta 
\pm 
\cdots 
=
\cdots 
\pm 
\cdots 
\int 
\int 
\int 
\rho 
\rho 
\rho 
\rho 
\rho 
\rho 
\rho ; 
\textcolor{lyred}{性质2}.
L
L
L
L
F dr
F dr
F dr
+
\cdots 
=
\cdots 
+
\cdots 
\int 
\int 
\int 
\rho 
\rho 
\rho 
\rho 
\rho 
\rho ; 
\textcolor{lyred}{性质3}.
L
L
F dr
F dr
-
\cdots 
=-
\cdots 
\int 
\int 
\rho 
\rho 
\rho 
\rho ; 
\textcolor{lyred}{性质4}.
\approx 
B
A
AB
dx
x
x
\cdots 
=
-
\int 
,
\approx 
B
A
AB
dy
y
y
\cdots 
=
-
\int 
. 
\vspace{4pt}
{\bfseries \textcolor{lyred}{四.计算法}}

\textcolor{lyred}{定理}.设光滑曲线
()
()
: x
x t
L
y
y t
=

=

, :t \alpha 
\beta 
\to 
,则
L
dr
F dr
F
dt
dt
\beta 
\alpha 


\cdots 
=
\cdots 




\int 
\int 
\rho 
\rho 
\rho 
\rho 
,即 
(
)
()
()
()
,
,
L
P x y dx
P x t
y t
x t dt
\beta 
\alpha 
'
=




\int 
\int 
,
(
)
()
()
()
,
,
L
Q x y dy
Q x t
y t
y t dt
\beta 
\alpha 
'
=




\int 
\int 
. 
\textcolor{lyred}{推论}.(1)若
()
:
L y
y x
=
,
:
x a
b
\to 
,则
(
)
(
)
,
,
L
P x y dx
Q x y dy
+
=
\int 
()
()
()
\{
\}
,
,
b
a
P x y x
Q x y x
y
x
dx
'
+








\int 
;(2)若
()
:
L x
x y
=
,
:
y c
d
\to 
,则 
(
)
(
)
()
()
()
\{
\}
,
,
,
,
d
L
c
P x y dx
Q x y dy
P x y
y x
y
Q x y
y
dy
'
+
=
+








\int 
\int 
. 
\textcolor{lyred}{例}.求
L
I
xydx
=\int 
,其中L 为
2y
x
=
上从(
)
1, 1
A
-
到(
)
1,1
B
的一段弧. 
解.
\approx 
\approx 
(
)
AO
OB
I
xydx
xydx
x
x dx
x xdx
x xdx
=
+
=
-
+
=
=
\int 
\int 
\int 
\int 
\int 
; 
或者,
I
y ydy
y dy
-
-
=
=
=
\int 
\int 
. 
\textcolor{lyred}{例}.设一质点在
(
)
,
M x y 处受到力F
\xi \rho 
的作用, F
\xi \rho 
的大小与M 到O 的距离成正比, 
方向指向O ,现在质点由(
)
,0
A a
沿椭圆
x
y
a
b
+
=按逆时针方向移动到
(
)
0,
B
b , 
求力F
\xi \rho 
所作的功. 
解.
cos
:
sin
x
a
t
L
y
b
t
=


=

, :0
t
\pi 
\to 
,
(
)
,
F
k x y
=-
\rho 
,
\approx 
\approx 
AB
AB
W
F dr
k
xdx
ydy
=
\cdots 
=-
+
=
\int 
\int 
\rho 
\rho 
(
)
(
)
(
)
cos
cos
sin
sin
sin cos
b
a
k a
td a
t
b
td b
t
k b
a
t
tdt
k
\pi 
\pi 


-
+
=-
-
=-
-




\int 
\int 
. 
\textcolor{lyred}{例}.求
L
I
xydx
=\int \nabla 
,其中L 为(
)
x
a
y
a
-
+
=
与x 轴所围成的区域位于第一象限 
部分的整个边界,取逆时针方向. 
解.
(
)
(
)
(
)
cos
sin
cos
a
L
L
I
xydx
xydx
x
dx
a
a
t a
t d a
a
t
\pi 
=
+
=
\cdots 
+
+
\cdots 
+
=
\int 
\int 
\int 
\int 
(
)
1 cos
sin
a
a
t
tdt
\pi 
\pi 
-
+
=-
\int 
. 
\textcolor{lyred}{例}.求
I
x dx
zy dy
x ydz
\Gamma 
=
+
-
\int 
,其中\Gamma 为
(
)
3,2,1
A
到O 的有向线段. 
解.
:
x
t
y
t
z
t
=


\Gamma 
=

=

,
: \to 
t
,故
I
t dt
=
=-
\int 
. 
\textcolor{lyred}{例}.求
(
)
(
)
(
)
I
z
y dx
x
z dy
x
y dz
\Gamma 
=
-
+
-
+
-
\int \nabla 
,其中
:
x
y
x
y
z

+
=
\Gamma -
+
=

,从z 轴 
正向看为顺时针方向. 
解.
cos
:
sin
cos
sin
x
t
y
t
z
t
t
=


\Gamma 
=

=
-
+

, : 2
t
\pi \to 
,故
(
)
4cos
I
t
dt
\pi 
\pi 
=
-
=-
\int 
. 
\vspace{4pt}
{\bfseries \textcolor{lyred}{五.曲线积分与路径无关的例子}}

\textcolor{lyred}{例}.求
L
I
xydx
x dy
=
+
\int 
,其中L 为(1)
y
x
=
上O 到(
)
1,1
B
的一段弧; 
(2)
x
y
=
上O 到(
)
1,1
B
的一段弧;(3)O 到(
)
1,0
A
,再到B 的折线. 
解.(1)
I
x x dx
x
dx
=
\cdots 
+
\cdots 
=
\int 
,(2)
(
)
I
y
y dy
y
dy
=
\cdots 
\cdots 
+
=
\int 
, 
(3)
(
)
(
)
,0
1,
OA
AB
I
P x
dx
Q
y dy
dy
=
+
=
+
=
+
=
\int 
\int 
\int 
\int 
\int 
. 
\vspace{4pt}
{\bfseries \textcolor{lyred}{六.两类曲线积分之间的联系}}

\textcolor{lyred}{定理}.设T
\rho 
为有向曲线L 的\textcolor{lyred}{切向量}(指向L 的正方向),则
(
)
L
L
F dr
F e
ds
\tau 
\cdots 
=
\cdots 
\int 
\int 
\rho 
\rho 
\rho 
\rho 
, 
\textcolor{lyred}{例}.设
:
L y
x
x
=
-
,(
)
(
)
0,0
1,1
\to 
,将
L
Pdx
Qdy
+
\int 
化为对弧长的曲线积分. 
解.
1,
1,
dy
x
T
dx
x
x


-


=
=





-


\rho 
,
(
)
,1
e
x
x
x
\tau =
-
-
\rho 
,故 
(
)(
)
(
)
,
,
L
L
Pdx
Qdy
x
x P x y
x Q x y
ds


+
=
-
+
-


\int 
\int 
. 
\textcolor{lyred}{例}.设
:
cos
L x
t
=
,
(
)
sin
y
t
t
\pi 
=
\le \le 
,取逆时针方向,证明: 
(
)
(
)
sin
cos 2
6 2
L
x
y dx
xy
dy
+
+
\le 
\int \nabla 
. 
证.由
(
)
(
)
,
,
L
L
L
L
Pdx
Qdy
P Q
e ds
P Q
e ds
P
Q ds
\tau 
\tau 
+
=
\cdots 
\le 
\cdots 
\le 
+
\int 
\int 
\int 
\int 
\rho 
\rho 
\nabla 
\nabla 
\nabla 
\nabla 
,得到 
(
)
(
)
sin
cos 2
6 2
L
L
x
y dx
xy
dy
ds
+
+
\le 
=
\int 
\int 
\nabla 
\nabla 
,证毕. 
\textcolor{lyred}{补充练习 }
1.求
L
xdy
ydx
I
x
y
-
=
+
\int 
,其中L 从
(
)
1,0
A -
沿
y
x
=-
-
到
(
)
1,0
B
,然后再沿直线到 
(
)
1,2
D -
的有向曲线. 
解. \approx 
cos
:
sin
x
t
AB
y
t
=


=

, :
t
\pi 
-
\to 
,
:
BD y
x
=-+,
:1
x
\to -,故 
4cos
sin
dt
dx
I
t
t
x
x
\pi 
\pi 
\pi 
\pi 
-
-
-
=
+
=
+
=
+
-
+
\int 
\int 
. 
2.求I
ydx
zdy
xdz
\Gamma 
=
+
+
\int \nabla 
,其中\Gamma 为
x
y
z
+
+
=被三个坐标面截下三角形的边界, 
从z 轴正向看为逆时针方向. 
解.设
1 :
x
x
y
x
z
=


\Gamma 
=-

=

,
:1
x
\to 
,由于在轮换x
y
z
x
\to 
\to 
\to 
下, I 不变,故 
(
)
I
ydx
zdy
xdz
ydx
x dx
\Gamma 
\Gamma 
=
+
+
=
=
-
=-
\int 
\int 
\int 
. 
\vspace{6pt}
{\large \bfseries \textcolor{lyred}{第11.3 节 格林公式及其应用}}

\vspace{4pt}
{\bfseries \textcolor{lyred}{一.Green 公式}}

\textcolor{lyred}{定理}.设分段光滑曲线L 为有界闭区域D 的正向边界, (
)
,
P x y 和
(
)
,
Q x y 在D 上 
有连续偏导数,则
L
D
Q
P
Pdx
Qdy
dxdy
x
y


\partial 
\partial 
+
=
-


\partial 
\partial 


\int 
\int \int 
\nabla 
. 
\textcolor{lyred}{注}. L
D
=\partial 
的正方向(诱导定向)为:当观察者在L 上沿这个方向行走时, D 总是 
位于他的左侧,即D 的外边界逆时针,内边界取顺时针. D
\partial 
的正向也记为D+
\partial 
. 
\textcolor{lyred}{推论}. 2
D
D
dxdy
xdy
ydx
\partial 
=
-
\int \int 
\int \nabla 
,或者
D
D
dxdy
xdy
ydx
\partial 
=
-
\int \int 
\int \nabla 
. 
\textcolor{lyred}{例}.求椭圆
x
y
a
a
+
=围成图形的面积. 
解.
(
)
(
)
cos
sin
sin
cos
D
D
dxdy
xdy
ydx
a
td b
t
b
td a
t
ab
\pi 
\pi 
\partial 
=
-
=
-
=
\int \int 
\int 
\int 
. 
\textcolor{lyred}{例}.求星形线
(
)
cos
sin
x
a
t
t
y
a
t
\pi 
=
\le \le 

=

围成图形的面积. 
解.
(
)
(
)
cos
sin
sin
cos
L
a
A
xdy
ydx
a
t
a
t
a
t
a
t
dt
\pi 
\pi 


'
'
=
-
=
\cdots 
-
\cdots 
=




\int 
\int 
\nabla 
. 
\textcolor{lyred}{例}.求
(
)
(
)
L
I
y
x dx
x
y dy
=
-
+
+
\int \nabla 
,其中
(
)
(
)
:
L
x
y
-
+
-
=
,逆时针. 
解.
(
)
3 1
D
D
Q
P
I
dxdy
dxdy
x
y
\pi 


\partial 
\partial 
=
-
=
-
=


\partial 
\partial 


\int \int 
\int \int 
. 
\textcolor{lyred}{例}.求
(
)
(
)
sin
cos
x
x
L
I
e
y
x
y dx
e
y
x dy
=
-
-
+
+
\int 
,其中L 为半圆周
y
x
x
=
-
上从 
(
)
0,0
O
到(
)
2,0
A
的一段有向弧段. 
解.
(
)
D
L AO
AO
I
dxdy
x dx
\pi 
+
=
-
=-
-
-
=-
-
\int 
\int 
\int \int 
\int 
\nabla 
. 
\textcolor{lyred}{例}.求
(
)
(
)
cos
1 2 sin
L
I
xy
y
x dx
y
x
x y
dy
=
-
+
-
+
\int 
,其中L 为
2x
y
\pi 
=
上从 
(
)
0,0
O
到
,1
A \pi 






的一段有向弧. 
解.设
,0
B \pi 






,则
1 2
L AB BO
AB
BO
I
y
y
dy
\pi 
\pi 
+
+


=
-
-
=-
-
+
=




\int 
\int 
\int 
\int 
\nabla 
. 
\textcolor{lyred}{例}.求
(
)
(
)
L
x
y dx
x
y dy
I
x
y
-
+
+
=
+
\int 
,其中L 为
cos
y
x
=
上从
,0
A
\pi 


-




到
,0
B \pi 






的 
一段有向弧. 
解.取
1 :
L
y
x
\pi 
=
-
上从B 到A 的一段,则
L L
L
L
I
+
=
-
=-
=
\int 
\int 
\int 
\nabla 
(
)
(
)
L
D
L
AB
AB
x
y dx
x
y dy
d
xdx
\pi 
\pi 
\sigma 
\pi 
\pi 
\pi 
\pi 
+
-




-
+
+


-
=-
-
=-
-
=-












\int 
\int 
\int 
\int \int 
\int 
\nabla 
. 
\textcolor{lyred}{例}.求
L
xdy
ydx
I
x
y
-
=
+
\int 
,其中L 从
,0
A
\pi 


-




沿
cos
y
x
\pi 


=
+




到
,0
B \pi 






. 
解.取
1L 为从
(
)
,0
M r
沿
x
y
r
+
=
逆时针到
(
)
,0
N
r
-
的有向弧,则 
L
L
L
L BM
L
NA
BM
NA
L
NM
NM
xdy
ydx
I
r
r
+
+
+
+


-
=
-
-
-
=-
=-
=-
-
=






\int 
\int 
\int 
\int 
\int 
\int 
\int 
\int 
\nabla 
\nabla 
D
dxdy
r
r
r
\pi 
\pi 
-
=-
\cdots 
=-
\int \int 
. 
\textcolor{lyred}{例}.求
L
xdy
ydx
I
x
y
-
=
+
\int \nabla 
,其中
:
x
y
L a
b
+
=,逆时针方向. 
解.取
:
r
C
x
y
r
+
=
,逆时针,则
(
)
r
r
r
L
C
C
C
xdy
ydx
I
r
\pi 
+-
-
-
=
-
=
=
\int 
\int 
\int 
\nabla 
\nabla 
\nabla 
. 
\textcolor{lyred}{例}.求
L
xdy
ydx
I
x
y
-
=
+
\int 
,其中L 为从
(
)
1,0
A -
沿
y
x
=-
-
到
(
)
1,0
B
,然后沿直线 
到
(
)
1,2
D -
的一段有向弧. 
解.取
: 4
C
x
y
r
+
=
,逆时针,则
C
L DA
DA
DA
xdy
ydx
I
r
+
-
=
-
=
-
=
\int 
\int 
\int 
\int 
\nabla 
\nabla 
x
y
r
dy
dxdy
r
y
\pi 
\pi 
\pi 
+
\le 
-
-
=
-
=
+
\int \int 
\int 
. 
\vspace{4pt}
{\bfseries \textcolor{lyred}{二.积分与路径无关的条件}}

\textcolor{lyred}{定理}.在区域G 内,积分值
L
Pdx
Qdy
+
\int 
与路径无关(只与起点终点有关) \Leftrightarrow 对于 
G 内的任意闭曲线C ,恒有
C
Pdx
Qdy
+
=
\int \nabla 
. 
\textcolor{lyred}{定义}.若平面区域G 内任意闭曲线所围部分均包含于G ,则称G 为\textcolor{lyred}{单连通区域},\textcolor{lyred}{ }
即无“洞”的区域,否则称为\textcolor{lyred}{复连通区域}. 
\textcolor{lyred}{定理}.设区域G 单连通,
(
)
,
P x y ,
(
)
,
Q x y 在G 内有连续偏导,则在G 内曲线积分 
L
Pdx
Qdy
+
\int 
与路径无关
Q
P
x
y
\partial 
\partial 
\Leftrightarrow 
=
\partial 
\partial 处处成立. 
\textcolor{lyred}{例}.求
(
)
y
y
L
I
ye dx
x
y
e dy
=
+
+
\int 
,其中L 为从O沿
x
y
=
到(
)
1,1
A
的有向弧. 
解. Q
P
x
y
\partial 
\partial 
=
\partial 
\partial , xOy 面单连通,故I 与路径无关,取(
)
1,0
B
,则 
(
)
(
)
(
)
,0
1,
y
OB
BA
I
P x
dx
Q
y dy
y
e dy
=
+
=
+
=
+
+
\int 
\int 
\int 
\int 
\int 
,此路不通; 
改取
(
)
0,1
C
,则
(
)
(
)
0,
,1
OC
CA
I
Q
y dy
P x
dx
e dx
e
=
+
=
+
=
+
\cdots 
=
\int 
\int 
\int 
\int 
\int 
. 
\textcolor{lyred}{例}.求
(
)
(
)
L
x
y dx
x
y dy
I
x
y
-
+
+
=
+
\int 
,其中L 为从
(
)
1,0
A -
沿
y
x
=
-
到(
)
1,0
B
的 
一段有向弧. 
解. Q
P
x
y
\partial 
\partial 
=
\partial 
\partial ,故I 与路径无关,取
1L 为从
(
)
1,0
A -
沿
y
x
=-
-
到(
)
1,0
B
的一段 
圆弧,则
(
)
(
)
L
L
I
x
y dx
x
y dy
\pi 
=
=
-
+
+
=
\int 
\int 
. 
\textcolor{lyred}{例}.求
(
)
L
xdy
ydx
I
x
y
-
=
-
\int 
,其中L 为
y
x
=
+
上从(
)
0,1
A
到(
)
3,2
B
的一段有向弧. 
解. Q
P
x
y
\partial 
\partial 
=
\partial 
\partial ,\{
\}
y
x
>
单连通,故
(
)
(
)
(
)
(
)
(
)
3,2
0,2
0,1
0,2
dx
I
x
-
=
+
=
+
=
-
-
\int 
\int 
\int 
. 
\vspace{4pt}
{\bfseries \textcolor{lyred}{三.全微分求积}}

{\bfseries \textcolor{lyred}{1.可积的条件}}

\textcolor{lyred}{定义}.若在平面区域G 内Pdx
Qdy
+
是某个(
)
,
u x y 的全微分,则称它在G 内\textcolor{lyred}{可积}, 
称(
)
,
u x y 为它的一个\textcolor{lyred}{原函数}. 
\textcolor{lyred}{定理}.设(
)
,
P x y ,
(
)
,
Q x y 均连续,若
L
Pdx
Qdy
+
\int 
在G 内与路径无关,则Pdx
Qdy
+
在G 内可积, (
)
(
)
(
)
,
,
,
x y
x
y
u x y
Pdx
Qdy
=
+
\int 
为它的一个原函数. 
\textcolor{lyred}{定理}.设区域G 单连通, (
)
,
P x y ,
(
)
,
Q x y 在G 内均有连续偏导数,则在G 内 
Pdx
Qdy
+
可积
Q
P
x
y
\partial 
\partial 
\Leftrightarrow 
=
\partial 
\partial 处处成立. 
\textcolor{lyred}{例}.验证
xdy
ydx
x
y
-
+
在右半平面内可积,并求它的一个原函数. 
解. Q
P
x
y
\partial 
\partial 
=
\partial 
\partial ,\{
\}
x >
单连通,故可积, (
)
(
)
(
)
,
1,0
,
x y
u x y
Pdx
Qdy
=
+
=
\int 
(
)
(
)
(
)
(
)
(
)
(
)
,0
,
1,0
,0
,0
,
arctan
x
x y
y
x
x
dx
xdy
y
P x
dx
Q x y dy
x
x
y
x
-\cdots 
+
=
+
=
+
+
\int 
\int 
\int 
\int 
. 
\textcolor{lyred}{定理}.设G 为单连通区域,
(
)
,
P x y ,
(
)
,
Q x y 在G 内有连续偏导数,则下面的四个 
命题等价: 
(1)曲线积分
L
Pdx
Qdy
+
\int 
在G 内与路径无关; 
(2)对于G 内任意闭曲线C ,均有
C
Pdx
Qdy
+
=
\int \nabla 
; 
(3) Q
P
x
y
\partial 
\partial 
=
\partial 
\partial 在G 内处处成立; 
(4)微分形式Pdx
Qdy
+
在G 内可积. 
{\bfseries \textcolor{lyred}{2.分项组合法}}

\textcolor{lyred}{例}.验证在xOy 面内,(
)
(
)
x
xy
dx
x y
y
dy
+
+
+
可积,求原函数. 
解.
(
)
(
)
Pdx
Qdy
x dx
xy dx
x ydy
y dy
dx
d
x y
dy
+
=
+
+
+
=
+
+
= 
(
)
d x
x y
y
+
+
,故(
)
,
u x y
x
x y
y
=
+
+
. 
\textcolor{lyred}{例}.验证在上半平面内,
x
xydx
dy
y


+
+




可积,求原函数. 
解.
ln
x
x y
Pdx
Qdy
xydx
dy
dy
d
y
y


+
=
+
+
=
+




,故(
)
,
ln
x y
u x y
y
=
+
. 
\vspace{4pt}
{\bfseries \textcolor{lyred}{四.曲线积分的基本定理}}

\textcolor{lyred}{定理}.设在G 内Pdx
Qdy
du
+
=
,则对G 内任意分段光滑有向弧段
\approx 
L
AB
=
,均有 
\approx 
\approx 
()
()
AB
AB
Pdx
Qdy
du
u B
u A
+
=
=
-
\int 
\int 
. 
\textcolor{lyred}{向量形式}:若
grad
F
u
=
\rho 
,则
\approx 
\approx 
()
()
AB
AB
F dr
du
u B
u A
\cdots 
=
=
-
\int 
\int 
\rho 
\rho 
. 
\textcolor{lyred}{注}.故当力场F
\xi \rho 
为保守场,即具有势函数u 时,它对在其中运动的质点作的功 
()
()
grad
L
L
L
W
F dr
u dr
du
u B
u A
=
\cdots 
=
\cdots 
=
=
-
\int 
\int 
\int 
\xi \rho 
\rho 
\rho 
,与路径无关. 
\textcolor{lyred}{例}.设右半平面内有一个力场,力的大小与点到原点的距离平方成正比,方向指向 
原点,证明:这个力场对在其中运动的质点所作的功与质点的运动路径无关,并求 
质点从(
)
1,1 到(
)
2,2 时力场对其所作的功. 
解.
(
)
(
)
(
)
3 2
,
,
grad
k
F
k
x
y
x y
k x x
y
y x
y
x
y
=-
+
=-
+
+
=-
+
\xi \rho 
,故 
L
W
F dr
=
\cdots 
\int 
\rho 
\rho 与路径无关,
(
)
(
)
(
)
2,2
1,1
14 2
k
W
x
y
k


=-
+
=-




;或者, 
Q
P
x
y
\partial 
\partial 
=
\partial 
\partial ,
x >
单连通,故
14 2
W
k x x
dx
k y
y dy
k
=-
+
-
+
=-
\int 
\int 
. 
\vspace{4pt}
{\bfseries \textcolor{lyred}{五.全微分方程}}

\textcolor{lyred}{定义}.一阶微分方程(
)
(
)
,
,
P x y dx
Q x y dy
+
=
称为\textcolor{lyred}{全微分方程},若存在可微函数 
(
)
,
u x y ,使得
(
)
(
)
,
,
du
P x y dx
Q x y dy
=
+
. 
\textcolor{lyred}{定理}.若(
)
(
)
(
)
,
,
,
P x y dx
Q x y dy
du x y
+
=
,则全微分方程 
(
)
(
)
,
,
P x y dx
Q x y dy
+
=
的通解为(
)
,
u x y
C
=
. 
\textcolor{lyred}{定理}.设在单连通区域内
(
)
,
P x y 和
(
)
,
Q x y 均有连续的偏导数,则微分方程 
(
)
(
)
,
,
P x y dx
Q x y dy
+
=
是全微分方程
Q
P
x
y
\partial 
\partial 
\Leftrightarrow 
=
\partial 
\partial . 
\textcolor{lyred}{例}.求(
)
(
)
x
y dx
x
y dy
-
-
-
=
通解. 
解一.
Q
P
x
y
\partial 
\partial 
=
=-
\partial 
\partial 
,又xOy 面单连通,故(
)
(
)
2x
y dx
x
y dy
-
-
-
可积, 
(
)
(
)
(
)
(
)
(
)
(
)
(
)
,
0,0
,
x y
y
x
u x y
x
y dx
x
y dy
x
dx
x
y
dy
=
-
-
-
=
-
+
-
-
=




\int 
\int 
\int 
x
xy
y
-
+
,故通解为
x
xy
y
C
-
+
=
,即
x
xy
y
C
-
+
=
; 
解二.(
)
(
)
x
y dx
x
y dy
x dx
ydx
xdy
ydy
-
-
-
=
-
-
+
= 
(
)
(
)
x dx
ydx
xdy
ydy
d
x
d xy
d
y
d
x
xy
y






-
+
+
=
-
+
=
-
+












, 
故通解为
x
xy
y
C
-
+
=
,即
x
xy
y
C
-
+
=
; 
解三.设(
)
(
)
2x
y dx
x
y dy
du
-
-
-
=
,则 
(
)
()
u
x
y
u
x
y dx
x
yx
C y
x
\partial 
=
-
\Rightarrow 
=
-
=
-
+
\partial 
\int 
,而由u
y
x
y
\partial 
=
-
\partial 
,得 
()
()
()
x
C
y
y
x
C
y
y
C y
y
C
'
'
-+
=
-
\Rightarrow 
=
\Rightarrow 
=
+
,故 
u
x
yx
y
C
=
-
+
+
,通解为
x
xy
y
C
-
+
=
,即
x
xy
y
C
-
+
=
. 
\textcolor{lyred}{例}.设
:
L x
y
x
y
+
+
+
=
,逆时针,证明:
cos
sin
L
x
y dy
y
x dx
\pi 
-
\le 
\int \nabla 
. 
证.左式
(
)
cos
sin
2 2
D
y
x
d
\pi 
\pi 
\sigma 
=
+
\le 
\cdots 
=
\int \int 
,证毕. 
\textcolor{lyred}{补充练习 }
1.求
(
)
(
)
y
L
I
x
xy
dx
e
x y dy
=
+
+
+
+
\int 
,其中
:
L y
x
=
-
,
:1
x
\to 
. 
解. Q
P
x
y
\partial 
\partial 
=
\partial 
\partial ,故
(
)
y
I
x
dx
e dy
e
=
+
+
=
-
\int 
\int 
. 
2.求
(
)
(
)
cos
2 sin
L
I
xy
y
x dx
x
y
x
x y
dy
=
-
+
-
+
\int 
,其中L 为从
,1
A \pi 






沿 
2x
y
\pi 
=
到
, 1
B \pi 


-




的一段有向弧. 
解.
D
L BA
BA
y
I
xdxdy
y
dx
\pi 
\pi 
\pi 
-
+


=
-
=
-
-
+
=-




\int 
\int 
\int \int 
\int 
\nabla 
. 
3.求
(
)
(
)
(
)
L
y
x dx
x
y dy
I
x
y
-
-
-
=
+
\int 
,其中L 从
(
)
1,0
A
沿
y
x
=
-
到
(
)
0,1
B
. 
解.
(
)
(
)
L BA
BA
AB
I
y
x dx
x
y dy
dx
+
=
-
=
-
-
-
=
\cdots 
=-
\int 
\int 
\int 
\int 
\nabla 
. 
4.求
(
)
L
I
x
y
xdx
ydy
=
+
+
\int \nabla 
,其中
(
)
:
L
x
y
-
+
=
,逆时针. 
解. Q
P
x
y
\partial 
\partial 
=
\partial 
\partial ,取
: 4
r
C
x
y
r
+
=
,逆时针,则
r
r
C
C
I
r
xdx
ydy
=
=
+
=
\int 
\int 
\nabla 
\nabla 
. 
5.设
()
f x >
有连续导数,且
()
f
=
,若在区域
x >
内,曲线积分 
()
()
ln
x
L
y
ye f x
dx
f x dy
x


-
-




\int 
与路径无关,求
()
f x . 
解.
()
()
()
x
Q
P
f
x
f x
e f
x
x
y
x
\partial 
\partial 
'
=
\Rightarrow 
-
=-
\partial 
\partial 
,即
()
y
f x
=
满足伯努利方程 
x
y
y
e y
x
'-
=-
,解得
x
x
x
y
xe
e
C
=
-
+
,又
()
f
=
,得
C =
. 
\vspace{6pt}
{\large \bfseries \textcolor{lyred}{第11.4 节 对面积的曲面积分}}

\vspace{4pt}
{\bfseries \textcolor{lyred}{一.第一类曲面积分的定义}}

\textcolor{lyred}{定义}.设\Sigma 为光滑曲面, (
)
, ,
f x y z 在\Sigma 上有界,将\Sigma 任意划分成n 片
iS
\Delta 
,直径为
i\lambda , 
面积为
iS
\Delta 
,作和
(
)
,
,
n
i
i
i
i
i
f
S
\psi \eta 
=
\Delta 
\sum 
,
(
)
,
,
i
i
i
iS
\psi \eta 
\forall 
\in \Delta 
,若在无限细分\Sigma 的过程中,
随着
\{\}
1max
i
i n
\lambda 
\lambda 
\le \le 
=
\to 
,该和趋向于一个只依赖(
)
, ,
f x y z 和\Sigma 的常数I ,则称I 为 
(
)
, ,
f x y z 在\Sigma 上\textcolor{lyred}{对面积的曲面积分},记为
(
)
, ,
f x y z dS
\Sigma \int \int 
,
(
)
, ,
f x y z 为\textcolor{lyred}{被积函数}, 
\Sigma 为\textcolor{lyred}{积分曲面}. 
\textcolor{lyred}{定理}.设
(
)
, ,
f x y z 在光滑曲面\Sigma 上连续,则
(
)
, ,
f x y z dS
\Sigma \int \int 
存在. 
\textcolor{lyred}{物理意义}.设\Sigma 为曲面型材料,密度为
(
)
, ,
x y z
\rho 
,连续,则
(
)
, ,
M
x y z dS
\rho 
\Sigma 
=\int \int 
. 
\vspace{4pt}
{\bfseries \textcolor{lyred}{二.第一类曲面积分的性质}}

\textcolor{lyred}{性质1}.
(
)
(
)
(
)
(
)
, ,
, ,
, ,
, ,
f x y z
g x y z
dS
f x y z dS
f x y z dS
\alpha 
\beta 
\alpha 
\beta 
\Sigma 
\Sigma 
\Sigma 
\pm 
=
\pm 




\int \int 
\int \int 
\int \int 
. 
\textcolor{lyred}{性质2}.
(
)
(
)
(
)
, ,
, ,
, ,
f x y z dS
f x y z dS
f x y z dS
\Sigma +\Sigma 
\Sigma 
\Sigma 
=
+
\int \int 
\int \int 
\int \int 
. 
\textcolor{lyred}{性质3}.
1 dS
S
\Sigma 
\cdots 
=
\int \int 
,一般地,
k dS
kS
\Sigma 
\cdots 
=
\int \int 
. 
\textcolor{lyred}{性质4}.若
(
)
(
)
, ,
, ,
f x y z
g x y z
\le 
,则
(
)
(
)
, ,
, ,
f x y z dS
g x y z dS
\Sigma 
\Sigma 
\le 
\int \int 
\int \int 
. 
\textcolor{lyred}{推论}.
(
)
(
)
, ,
, ,
f x y z dS
f x y z dS
\Sigma 
\Sigma 
\le 
\int \int 
\int \int 
. 
\textcolor{lyred}{注}.也有估值定理和积分中值定理. 
\vspace{4pt}
{\bfseries \textcolor{lyred}{三.对称性}}

\textcolor{lyred}{注}.参照三重积分.\textcolor{lyred}{ }
\textcolor{lyred}{例}.求
(
)
I
x
xz dS
\Sigma 
=
+
\int \int 
,其中
(
)
:
2 0
x
y
z
\Sigma 
+
=
\le 
\le 
. 
解.
(
)
8 2
I
x dS
x
y
dS
dS
\pi 
\Sigma 
\Sigma 
\Sigma 
=
=
+
=
=
\int \int 
\int \int 
\int \int 
. 
\textcolor{lyred}{例}.求
(
)
I
x
y dS
\Sigma 
=
+
\int \int 
,其中
:
x
y
z
\Sigma 
+
+
=. 
解.
(
)
I
y dS
x
y
z dS
dS
\Sigma 
\Sigma 
\Sigma 
=
=
+
+
=
=
\cdots \cdots 
=
\int \int 
\int \int 
\int \int 



. 
\textcolor{lyred}{例}.求
(
)
I
x
y
z
dS
\Sigma 
=
-
+
+
\int \int 
,其中
:
x
y
z
\Sigma 
+
+
=. 
解.
(
)
I
x
y
z
dS
x dS
dS
dS
\pi 
\Sigma 
\Sigma 
\Sigma 
\Sigma 
=
+
+
+
=
+
=
=
\int \int 
\int \int 
\int \int 
\int \int 




. 
\textcolor{lyred}{例}.求
(
)
I
x
y
z
dS
\Sigma 
=
+
+
\int \int 
,其中
:
x
y
z
x
\Sigma 
+
+
=
. 
解.
(
)
(
)
I
x
y
z
xy
yz
zx dS
x
xy
yz
zx dS
\Sigma 
\Sigma 
=
+
+
+
+
+
=
+
+
+
=
\int \int 
\int \int 


(
)
1 1
xdS
x
dS
dS
\pi 
\Sigma 
\Sigma 
\Sigma 
=
-+
=
=
\int \int 
\int \int 
\int \int 



. 
\textcolor{lyred}{例}.求
(
)
I
x
y
z
dS
\Sigma 
=
+
+
\int \int 
,其中
:
x
y
z
y
\Sigma 
+
+
=
. 
解.
(
)
(
)
I
x
y
dS
x
y
z
dS
ydS
dS
\pi 
\Sigma 
\Sigma 
\Sigma 
\Sigma 
=
+
=
+
+
=
=
=
\int \int 
\int \int 
\int \int 
\int \int 




. 
\textcolor{lyred}{例}.求
(
)
I
x
y
z
dS
\Sigma 
=
+
+
-
\int \int 
,其中
(
)
(
)
(
)
:
x
y
z
\Sigma 
-
+
-
+
-
=. 
解.
(
)(
)(
)
I
x
y
z
dS
\Sigma 
=
-
+
-
+
-
+
=




\int \int 
(
)
(
)
(
)
(
)
x
y
z
dS
dS
dS
\pi 
\Sigma 
\Sigma 
\Sigma 


-
+
-
+
-
+
=
+
=
=


\int \int 
\int \int 
\int \int 



. 
\vspace{4pt}
{\bfseries \textcolor{lyred}{四.计算法}}

\textcolor{lyred}{定理(投影法)}.设
(
)
, ,
f x y z 在光滑曲面\Sigma 上连续,记
(
)
, ,
I
f x y z dS
\Sigma 
=\int \int 
,则 
(1)当
(
)
:
,
z
z x y
\Sigma 
=
,(
)
,
xy
x y
D
\in 
时,
(
)
, ,
,
xy
x
y
D
I
f
x y z x y
z
z dxdy
=
+
+




\int \int 
; 
(2)当
(
)
:
,
y
y x z
\Sigma 
=
,(
)
,
xz
x z
D
\in 
时,
(
)
,
,
,
xz
x
z
D
I
f
x y x z
z
y
y dxdz
=
+
+




\int \int 
; 
(3)当
(
)
:
,
x
x y z
\Sigma 
=
,(
)
,
yz
y z
D
\in 
时,
(
)
,
, ,
yz
y
z
D
I
f
x y z
y z
x
x dydz
=
+
+




\int \int 
. 
\textcolor{lyred}{例}.求I
xyzdS
\Sigma 
=\int \int 
,其中\Sigma 为平面
x
y
z
+
+
=与三个坐标面围成四面体的表面. 
解.
(
)
1,
0,
x y
x
y
I
xyz dS
xy
x
y
dxdy
\Sigma 
\Sigma 
\Sigma 
\Sigma 
\Sigma 
+\le 
\ge 
\ge 
=
+
+
+
=
\cdots 
=
-
-
\cdots 
=
\int \int 
\int \int 
\int \int 
\int \int 
\int \int 
\int \int 
(
)
x
dx
xy
x
y dy
-
-
-
=
\int 
\int 
. 
\textcolor{lyred}{例}.求
(
)
I
xyz y z
z x
x y
dS
\Sigma 
=
+
+
\int \int 
,其中
(
)
:
, ,
x
y
z
a
x y z
\sum 
+
+
=
\ge 
位于. 
解.
a
z
a
x
y
dS
d
a
x
y
\sigma 
=
-
-
\Rightarrow 
=
-
-
,故
I
xyz x y dS
\Sigma 
=
\cdots 
=
\int \int 
xy
D
a
a
x y
a
x
y
d
a
x
y
\sigma 
-
-
\cdots 
=
-
-
\int \int 
. 
\textcolor{lyred}{例}.求
(
)
I
x
y
z dS
\Sigma 
=
+
+
\int \int 
,其中\sum 为
y
z
+
=
被
x
y
+
=
截下部分. 
解.
z
y
dS
d\sigma 
=
-
\Rightarrow 
=
,故
(
)
2 25
125 2
x
y
I
y
z dS
d\sigma 
\pi 
\Sigma 
+
\le 
=
+
=
\cdots 
=
\int \int 
\int \int 
. 
\textcolor{lyred}{例}.求
(
)
I
xy
yz
zx dS
\Sigma 
=
+
+
\int \int 
,其中\Sigma 为
z
x
y
=
+
被
x
y
ax
+
=
截下部分. 
解.
2 cos
64 2
cos
xy
a
D
I
zxdS
x x
y
d
d
d
a
\pi 
\theta 
\pi 
\sigma 
\theta 
\rho 
\theta \rho 
\Sigma 
-
=
=
+
\cdots 
=
=
\int \int 
\int \int 
\int 
\int 
. 
\textcolor{lyred}{例}.求
dS
I
x
y
z
\Sigma 
=
+
+
\int \int 
,其中\Sigma 为
(
)
2 0
x
y
R
z
H
+
=
\le 
\le 
. 
解.设
1 : x
R
y
\Sigma 
=
-
,则
x
x
R
dS
dydz
dydz
y
z
R
y


\partial 
\partial 


=
+
+
=




\partial 
\partial 


-


, 
yz
R
H
D
R
dS
R
Rdy
dz
I
dydz
x
y
z
R
z
R
z
R
y
R
y
\Sigma 
-
=
=
\cdots 
=
\cdots 
=
+
+
+
+
-
-
\int \int 
\int \int 
\int 
\int 
2 arctan H
R
\pi 
. 
\textcolor{lyred}{例}.求I
xdS
\Sigma 
=\int \int 
,其中\Sigma 为
x
y
+
=,
z
x
=
+
,
z =
围成立体的表面. 
解.设
1 :
z
x
\Sigma 
=
+
,
2 :
z
\Sigma 
=
,
3 :
y
x
\Sigma 
=
-
,则
I
\Sigma 
\Sigma 
\Sigma 
=
+
+
=
\int \int 
\int \int 
\int \int 
xy
xy
xz
x
D
D
D
x
x
dxdy
x dxdy
x
dxdz
dx
dz
x
x
\pi 
+
-
\cdots 
+
\cdots 
+
\cdots 
=
+
+
=
-
-
\int \int 
\int \int 
\int \int 
\int 
\int 
. 
\vspace{4pt}
{\bfseries \textcolor{lyred}{五.物理应用}}

设\Sigma 为曲面型材料,密度为
(
)
, ,
x y z
\mu 
,则 
1.质心坐标为
x dS
x
dS
\mu 
\mu 
\Sigma 
\Sigma 
=\int \int 
\int \int 
,
y dS
y
dS
\mu 
\mu 
\Sigma 
\Sigma 
=\int \int 
\int \int 
,
z dS
z
dS
\mu 
\mu 
\Sigma 
\Sigma 
=\int \int 
\int \int 
; 
2.对l 的转动惯量
lI
d
dS
\mu 
\Sigma 
=\int \int 
,其中(
)
, ,
d x y z 为(
)
, ,
x y z 到l 的距离; 
3.对位于(
)
,
,
x
y
z
处的单位质点的引力 
(
)
(
)
(
)
,
,
x
x
y
y
z
z
F
G
dS G
dS G
dS
r
r
r
\mu 
\mu 
\mu 
\Sigma 
\Sigma 
\Sigma 


-
-
-
=



\int \int 
\int \int 
\int \int 
\rho 
,其中 
(
)
(
)
(
)
r
x
x
y
y
z
z
=
-
+
-
+
-
. 
\textcolor{lyred}{例}.求密度为\mu 的均匀半球形壳
: z
a
x
y
\Sigma 
=
-
-
对位于原点处的单位质点的 
引力. 
解.由对称性,
x
y
F
F
=
=
,
(
)
(
)
z
G
z
G
F
dS
zdS
a
x
y
z
\mu 
\mu 
\Sigma 
\Sigma 
-
=
=
=
+
+
\int \int 
\int \int 
x
y
a
G
a
a
x
y
dxdy
G
a
a
x
y
\mu 
\mu \pi 
+
\le 
-
-
\cdots 
=
-
-
\int \int 
. 
\textcolor{lyred}{补充练习 }
1.求
dS
I
z
\Sigma 
=\int \int 
,其中\Sigma 是
x
y
z
+
+
=
被
z =截下的顶部. 
解.
4 ln 2
x
y
d
I
dxdy
d
a
x
y
x
y
\pi 
\rho \rho 
\theta 
\pi 
\rho 
+
\le 
=
\cdots 
=
=
-
-
-
-
-
\int \int 
\int 
\int 
. 
2.求I
xyz dS
\Sigma 
=\int \int 
,其中
(
)
:
z
x
y
z
\Sigma 
=
+
\le 
\le 
. 
解.
(
)
1,
0,
x
y
x
y
I
xyz dS
xy x
y
x
y dxdy
\Sigma 
+
\le 
\ge 
\ge 
=
\cdots 
=
+
+
+
=
\int \int 
\int \int 
125 5
cos sin
d
d
\pi 
\theta \rho 
\theta 
\theta \rho 
\rho 
\rho \rho 
-
\cdots 
+
\cdots 
=
\int 
\int 
. 
\vspace{6pt}
{\large \bfseries \textcolor{lyred}{第11.5 节 对坐标的曲面积分}}

\vspace{4pt}
{\bfseries \textcolor{lyred}{一.第二类曲面积分的定义}}

我们常见的曲面大都是\textcolor{lyred}{双侧的},选定一侧的双侧曲面称为\textcolor{lyred}{有向曲面}. 
设\Sigma 为光滑有向曲面,
(
)
, ,
R x y z 在\Sigma 上有界,将\Sigma 任意分成n 小片
iS
\Delta 
,直径为
i\lambda , 
作和
(
)(
)
,
,
n
i
i
i
i
xy
i
R
S
\psi \eta 
=
\Delta 
\sum 
, (
)
,
,
i
i
i
iS
\psi \eta 
\forall 
\in \Delta 
,其中(
)
i
xy
S
\Delta 
为
iS
\Delta 
在xOy 平面上的 
投影,若在无限细分\Sigma 的过程中,随着
\{\}
1max
i
i n
\lambda 
\lambda 
\le \le 
=
\to 
,该和趋向于一个只依赖 
于
(
)
, ,
R x y z 和\Sigma 的常数I ,则称I 为
(
)
, ,
R x y z 在\Sigma 上\textcolor{lyred}{对坐标x 与y 的曲面积分}, 
记为
(
)
, ,
R x y z dxdy
\Sigma \int \int 
;称
(
)
, ,
f x y z 为\textcolor{lyred}{被积函数},\Sigma 为\textcolor{lyred}{积分曲面}. 
类似地,
(
)
(
)(
)
, ,
lim
,
,
n
i
i
i
i
yz
i
P x y z dydz
P
S
\lambda 
\psi \eta 
\to 
=
\Sigma 
=
\Delta 
\sum 
\int \int 
; 
(
)
(
)(
)
, ,
lim
,
,
n
i
i
i
i
zx
i
Q x y z dzdx
Q
S
\lambda 
\psi \eta 
\to 
=
\Sigma 
=
\Delta 
\sum 
\int \int 
. 
设
(
)
,
,
A
P Q R
=
\rho 
为\Sigma 上的有界向量值函数(场),则A
\rho 
在\Sigma 上\textcolor{lyred}{对坐标的曲面积分}为 
A dS
Pdydz
Qdzdx
Rdxdy
\Sigma 
\Sigma 
\cdots 
=
+
+
\int \int 
\int \int 
\rho 
\rho 
,其中
(
)
,
,
dS
dydz dzdx dxdy
=
\rho 
为\textcolor{lyred}{有向曲面元素}. 
\textcolor{lyred}{定理}.光滑有向曲面上的连续函数对坐标的曲面积分存在. 
\vspace{4pt}
{\bfseries \textcolor{lyred}{二.两类曲面积分的联系}}

\textcolor{lyred}{定理}.设光滑有向曲面\Sigma 的单位法向量
(
)
cos ,cos
,cos
ne
\alpha 
\beta 
\gamma 
=
\rho 
,则 
cos
Pdydz
P
dS
\alpha 
\Sigma 
\Sigma 
=
\cdots 
\int \int 
\int \int 
,
cos
Qdzdx
Q
dS
\beta 
\Sigma 
\Sigma 
=
\cdots 
\int \int 
\int \int 
,
cos
Rdxdy
R
dS
\gamma 
\Sigma 
\Sigma 
=
\int \int 
\int \int 
,即 
(
)
(
)
cos
cos
cos
n
A dS
A e
dS
P
Q
R
dS
\alpha 
\beta 
\gamma 
\Sigma 
\Sigma 
\Sigma 
\cdots 
=
\cdots 
=
+
+
\int \int 
\int \int 
\int \int 
\rho 
\rho 
\rho \rho 
. 
\textcolor{lyred}{例}.设\Sigma 为圆柱体(
)
(
)
(
)
4 1
x
x
y
y
z
-
+
-
\le 
\le 
\le 
的外表面,证明: 
(
)
(
)
cos
sin 2
16 3
x
y dydz
xy
dzdx
dxdy
\pi 
\Sigma 
+
+
+
\le 
\int \int 
. 
证.
(
)
n
n
A dS
A e
dS
A e dS
A dS
dS
\Sigma 
\Sigma 
\Sigma 
\Sigma 
\Sigma 
\cdots 
=
\cdots 
\le 
\cdots 
\le 
\le 
\int \int 
\int \int 
\int \int 
\int \int 
\int \int 
\rho 
\rho 
\rho 
\rho 
\rho 
\rho 
\rho 
,证毕. 
\vspace{4pt}
{\bfseries \textcolor{lyred}{三.流过曲面一侧的流量}}

设
(
)
,
,
A
P Q R
=
\rho 
为一稳定不可压缩流体的流速场,\Sigma 为一片有向曲面,在\Sigma 上任取 
dS ,则单位时间内流过其指定一侧的流体构成了一个底面积为dS ,斜高为A
\rho 
的 
斜柱体,体积
(
)
n
d
A e
dS
A dS
\Phi =
\cdots 
=
\cdots 
\rho 
\rho 
\rho 
\rho 
,故单位时间内流过\Sigma 指定一侧的流量总量 
A dS
\Sigma 
\Phi =
\cdots 
\int \int 
\rho 
\rho 
,也称为\textcolor{lyred}{通量}. 
\vspace{4pt}
{\bfseries \textcolor{lyred}{四.第二类曲面积分的性质}}

\textcolor{lyred}{性质1}.
(
)
A
B
dS
A dS
B dS
\alpha 
\beta 
\alpha 
\beta 
\Sigma 
\Sigma 
\Sigma 
\pm 
\cdots 
=
\cdots 
\pm 
\cdots 
\int \int 
\int \int 
\int \int 
\rho 
\rho 
\rho 
\rho 
\rho 
\rho 
; 
\textcolor{lyred}{性质2}.
A dS
A dS
A dS
\Sigma +\Sigma 
\Sigma 
\Sigma 
\cdots 
=
\cdots 
+
\cdots 
\int \int 
\int \int 
\int \int 
\rho 
\rho 
\rho 
\rho 
\rho 
\rho 
; 
\textcolor{lyred}{性质3}.
A dS
A dS
-\Sigma 
\Sigma 
\cdots 
=-
\cdots 
\int \int 
\int \int 
\rho 
\rho 
\rho 
\rho 
. 
\vspace{4pt}
{\bfseries \textcolor{lyred}{五.计算法}}

\textcolor{lyred}{定理(投影法)}.设\Sigma 为光滑有向曲面,
,
,
P Q R 为\Sigma 上的有界函数. 
(1)设
(
)
:
,
z
z x y
\Sigma 
=
,上侧,则
(
)
, ,
,
xy
D
Rdxdy
R x y z x y
dxdy
\Sigma 
=




\int \int 
\int \int 
; 
(2)设
(
)
:
,
y
y x z
\Sigma 
=
,右侧,则
(
)
,
,
,
xz
D
Qdzdx
Q x y x z
z dzdx
\Sigma 
=




\int \int 
\int \int 
; 
(3)设
(
)
:
,
x
x y z
\Sigma 
=
,前侧,则
(
)
,
, ,
yz
D
Pdydz
P x y z
y z dydz
\Sigma 
=




\int \int 
\int \int 
. 
\textcolor{lyred}{定理(合一公式)}.设I
Pdydz
Qdzdx
Rdxdy
\Sigma 
=
+
+
\int \int 
,则 
(1)当
(
)
:
,
z
z x y
\Sigma 
=
时,
z
z
I
P
Q
R
dxdy
x
y
\Sigma 




\partial 
\partial 


=
\cdots -
+
\cdots -
+
\cdots 






\partial 
\partial 






\int \int 
; 
(2)当
(
)
:
,
y
y x z
\Sigma 
=
时,
y
y
I
P
Q
R
dzdx
x
z
\Sigma 

\partial 
\partial 





=
\cdots -
+
\cdots +
\cdots -






\partial 
\partial 






\int \int 
; 
(3)当
(
)
:
,
x
x y z
\Sigma 
=
时,
x
x
I
P
Q
R
dydz
y
z
\Sigma 




\partial 
\partial 


=
\cdots +
\cdots -
+
\cdots -






\partial 
\partial 






\int \int 
. 
\textcolor{lyred}{例}.设\Sigma 为
[
][
][
]
0,
0,
0,
a
b
c
\Omega =
\times 
\times 
的外侧表面,求
I
x dydz
y dzdx
z dxdy
\Sigma 
=
+
+
\int \int 
. 
解.
xy
xy
D
D
z dxdy
c dxdy
dxdy
c ab
\Sigma 
\Sigma 
\Sigma 
=
+
=
-
=
\int \int 
\int \int 
\int \int 
\int \int 
\int \int 

,同理, 
y dzdx
b ac
\Sigma 
=
\int \int 
,
x dydz
a bc
\Sigma 
=
\int \int 
,故
(
)
I
a
b
c abc
=
+
+
. 
\textcolor{lyred}{例}.求
ze
I
dxdy
x
y
\Sigma 
=
+
\int \int 
,其中\Sigma 为
z
x
y
=
+
,
z =,
z =
围成立体的外表面. 
解.记\Sigma 的上,下,侧面为
\Sigma ,
\Sigma ,
\Sigma ,则
I
\Sigma 
\Sigma 
\Sigma 
=
+
+
=
\int \int 
\int \int 
\int \int 
x
y
x
y
x
y
x
y
e
e
e
dxdy
dxdy
dxdy
e
x
y
x
y
x
y
\pi 
+
+
\le 
+
\le 
\le 
+
\le 
-
-
=
+
+
+
\int \int 
\int \int 
\int \int 
. 
\textcolor{lyred}{例}.求
(
)
I
x
xyz dxdy
\Sigma 
=
+
\int \int 
,其中\Sigma 为
(
)
0,
x
y
z
x
y
+
+
=
\ge 
\ge 
外侧. 
解.
(
)
(
)
xy
xy
D
D
I
x
xy
x
y
dxdy
x
xy
x
y
dxdy
\Sigma 
\Sigma 
=
+
=
+
-
-
-
-
-
-
=
\int \int 
\int \int 
\int \int 
\int \int 
cos
sin
xy
D
xy
x
y dxdy
d
d
\pi 
\theta \rho 
\theta \rho 
\theta 
\rho 
\rho \rho 
-
-
=
\cdots 
\cdots 
-
\cdots 
=
\int \int 
\int 
\int 
. 
\textcolor{lyred}{例}.求
xdydz
z dxdy
I
x
y
z
\Sigma 
+
=
+
+
\int \int 
,其中\Sigma 为
\{
\}
1,
x
y
z
\Omega =
+
\le 
\le 
的外表面. 
解.记\Sigma 的上下前后面为
\Sigma ,
\Sigma ,
\Sigma ,
\Sigma ,则
I
\Sigma 
\Sigma 
\Sigma 
\Sigma 
=
+
+
+
=
\int \int 
\int \int 
\int \int 
\int \int 
(
)
(
)
yz
yz
D
D
x
y
x
y
dxdy
y
y
dxdy
dydz
dydz
x
y
z
z
x
y
+
\le 
+
\le 
-
-
-
-
\cdots 
-
+
-
=
+
+
+
+
+
+-
\int \int 
\int \int 
\int \int 
\int \int 
yz
D
y
y
dydz
dy
dz
z
z
\pi 
-
-
-
-
=
=
+
+
\int \int 
\int 
\int 
. 
\textcolor{lyred}{例}.求
(
)
(
)
I
x y
x
dydz
z
dxdy
\Sigma 
=
-
+
+
+
\int \int 
,其中\Sigma 为
x
y
+
=
被
z =
,
y
z
+
=
截下部分的外侧. 
解.设
1 :
x
y
\Sigma 
=
-
前侧,则 
(
)
(
)
yz
y
D
I
x dydz
y
dydz
dy
y dz
\pi 
-
\Sigma 
-
=
-
=
-
-
=-
-
=-
\int \int 
\int \int 
\int 
\int 
. 
\textcolor{lyred}{例}.求I
xdydz
zdxdy
\Sigma 
=
-
\int \int 
,其中
(
)(
)
:
z
x
y
z
\Sigma 
=
+
\le 
\le 
下侧. 
解.
(
)(
)
(
)(
)
,0,
,
,1
,0,
,
,1
x
y
I
x
z
z
z
dxdy
x
z
x
y
dxdy
\Sigma 
\Sigma 
=
-
\cdots -
-
=
-
\cdots -
-
=
\int \int 
\int \int 
(
)
(
)
(
)
x
y
x
y
x
z dxdy
x
x
y
d
x
y
d
\sigma 
\sigma 
\pi 
\Sigma 
+
\le 
+
\le 


-
-
=
+
+
=
+
=




\int \int 
\int \int 
\int \int 
. 
\textcolor{lyred}{例}.求
(
)
(
)
I
y
z dydz
x
y dxdy
\Sigma 
=
-
+
-
\int \int 
,其中\Sigma 为
z
x
x
y
=
-
-
被
x
y
x
+
=
截下部分的上侧. 
解.
x
x
y
x
zz
x
x
y
z
x
z
y
zz
z
+
=

-
+
+
=
\Rightarrow 
\Rightarrow 
=

+
=

,
y
y
z
z
-
=
,故 
(
)
(
)
(
)
,0,
,
,1
x
y
x
I
y
z
x
y
dxdy
y
z
x
y
dxdy
z
z
z
\Sigma 
\Sigma 
-
-
-




=
-
-
\cdots -
-
=
-
+
-
=








\int \int 
\int \int 
(
)
(
)
xy
D
x
y
x
x
y
x
x
y
x
y
d
d
x
x
y
\sigma 
\sigma 
\pi 
+
\le 


-
-
-
-
+
-
=
\cdots 
=



-
-



\int \int 
\int \int 
. 
\textcolor{lyred}{补充练习} 
1.求
I
zdxdy
xdydz
y dzdx
\Sigma 
=
+
+
\int \int 
,其中
(
)
:
1 0
x
y
z
\Sigma 
+
=
\le 
\le 
外侧. 
解.设
1 :
x
y
\Sigma 
=
-
前侧,则 
yz
D
I
xdydz
y dydz
dy
y dz
\pi 
\Sigma 
-
=
=
-
=
-
=
\int \int 
\int \int 
\int 
\int 
. 
2.求
(
)
I
y
z dydz
z dxdy
\Sigma 
=
+
+
\int \int 
,其中
(
)
:
z
x
y
z
\Sigma 
=
+
\le 
\le 
下侧. 
解.
(
)
,0,
,
,1
x
y
xy
xz
I
y
z
z
dxdy
z
dxdy
x
y
x
y
x
y
\Sigma 
\Sigma 




-
-
+




=
+
\cdots 
=
-
=




+
+
+




\int \int 
\int \int 
(
)
x
y
x
y
dxdy
d
d
\pi 
\theta \rho 
\rho \rho 
\pi 
\le 
+
\le 
-
+
=-
\cdots 
=-
\int \int 
\int 
\int 
. 
3.求
z
z
z
I
xe dydz
ye dzdx
e dxdy
\Sigma 
=
+
-
\int \int 
,其中
(
)
:
1 0
z
x
y
z
\Sigma 
=
+
-
\le 
\le 
下侧. 
解.
(
)
,
,
,
,1
z
z
z
x
y
x
y
I
xe
ye
e
dxdy
x
y
x
y
\le 
+
\le 


-
-


=-
-
\cdots 
=


+
+


\int \int 
(
)
(
)
(
)
x
y
x
y
e
x
y
d
d
e
d
e
\pi 
\rho 
\sigma 
\theta 
\rho 
\rho \rho 
\pi 
+
-
-
\le 
+
\le 
+
+
=
+
\cdots 
=
-
\int \int 
\int 
\int 
. 
\vspace{6pt}
{\large \bfseries \textcolor{lyred}{第11.6 节 高斯公式}}

\vspace{4pt}
{\bfseries \textcolor{lyred}{一.Gauss 公式}}

\textcolor{lyred}{定义}.设向量场
(
)
,
,
A
P Q R
=
\rho 
,记div
P
Q
R
A
x
y
z
\partial 
\partial 
\partial 
=
+
+
\partial 
\partial 
\partial 
\rho 
,称为A
\rho 
的\textcolor{lyred}{散度}.\textcolor{lyred}{ }
\textcolor{lyred}{定理}.设光滑有向闭曲面\Sigma 为闭区域\Omega 的外侧边界,
(
)
,
,
A
P Q R
=
\rho 
为\Omega 上向量场, 
其中
(
)
, ,
P x y z ,
(
)
, ,
Q x y z ,
(
)
, ,
R x y z 均在\Omega 上有连续偏导数,则 
div
A dS
A dv
\Sigma 
\Omega 
\cdots 
=
\cdots 
\int \int 
\int \int \int 
\rho 
\rho 
\rho 

. 
\textcolor{lyred}{推论}.
xdydz
ydzdx
zdxdy
\partial \Omega 
\Omega =
+
+
\int \int 
. 
\textcolor{lyred}{例}.求
(
)
(
)
I
y
x dxdy
x y
z dydz
\Sigma 
=
-
+
-
\int \int 
,其中\Sigma 为
x
y
+
=,
z =
,
z =
所围成 
立体的内表面. 
解.
(
)
(
)
(
)
div
,0,
I
x y
z
y
x dv
z
y dv
zdv
dv
\pi 
\Omega 
\Omega 
\Omega 
\Omega 
=-
-
-
=
-
=
=
=
\int \int \int 
\int \int \int 
\int \int \int 
\int \int \int 
. 
\textcolor{lyred}{例}.设(
)
()
, ,
u x y z
C
\in 
\Omega ,证明:
u
u
u
u
dS
dv
n
x
y
z
\partial \Omega 
\Omega 


\partial 
\partial 
\partial 
\partial 
=
+
+


\partial 
\partial 
\partial 
\partial 


\int \int 
\int \int \int 

. 
证.
(
)
grad
,
,
n
u
u
u
u
dS
u e
dS
dS
n
x
y
z
\partial \Omega 
\partial \Omega 
\partial \Omega 


\partial 
\partial 
\partial 
\partial 
=
\cdots 
=
\cdots 


\partial 
\partial 
\partial 
\partial 


\int \int 
\int \int 
\int \int 
\rho 
\rho 



,即得,证毕. 
\textcolor{lyred}{例}.求
(
)
xdydz
ydzdx
zdxdy
I
x
y
z
\Sigma 
+
+
=
+
+
\int \int 
,其中
:
x
y
z
a
b
c
\Sigma 
+
+
=,外侧. 
解.取
1 : x
y
z
r
\Sigma 
+
+
=
,外侧,
\{
\}
\ x
y
z
r
'
\Omega =\Omega 
+
+
\le 
,则 
(
)
xdydz
ydzdx
zdxdy
I
dv
r
r
r
\pi 
\pi 
'
\Sigma +-\Sigma 
-\Sigma 
\Omega 
\Sigma 
+
+
=
-
=
\cdots 
+
=
\cdots \cdots 
=
\int \int 
\int \int 
\int \int \int 
\int \int 



. 
\textcolor{lyred}{例}.求
I
x dydz
y dzdx
z dxdy
\Sigma 
=
+
+
\int \int 
,其中
(
)
:
x
y
z
z
h
\Sigma 
+
=
\le 
\le 
,下侧. 
解.取
(
)
1 : z
h x
y
h
\Sigma 
=
+
\le 
,上侧,则
I
\Sigma +\Sigma 
\Sigma 
=
-
=
\int \int 
\int \int 

(
)
h
x
y
h
x
y
z
h
x
y
z dv
h dxdy
dz
zdxdy
h
\pi 
\pi 
\Omega 
+
\le 
+
\le 
+
+
-
=
-
=-
\int \int \int 
\int \int 
\int 
\int \int 
. 
\textcolor{lyred}{例}.求
(
)
(
)
(
)
2 1
I
x
y
dydz
z
y
dzdx
yz
x
dxdy
\Sigma 
=
+
+
+
-
-
+
+
\int \int 
,其中 
(
)
:
y
x
z
y
\Sigma 
-=
+
\le 
\le 
,外侧. 
解.取
(
)
1 :
y
x
z
\Sigma 
=
+
\le 
,右侧,则
I
\Sigma +\Sigma 
\Sigma 
=
-
=
\int \int 
\int \int 

(
)
2 1 3
x
z
x
z
y
dv
dzdx
dy
dzdx
\pi 
\pi 
\Omega 
+
\le 
+
\le -
-
-
=
+
=
\int \int \int 
\int \int 
\int 
\int \int 
. 
\textcolor{lyred}{例}.求
z
z
z
I
xe dydz
ye dzdx
e dxdy
\Sigma 
=
+
-
\int \int 
,其中
(
)
:
1 0
z
x
y
z
\Sigma 
=
+
-
\le 
\le 
,下侧. 
解.取
(
)
1 :
z
x
y
\Sigma 
=
+
\le 
,上侧,
(
)
2 :
z
x
y
\Sigma 
=
+
\le 
,下侧,则 
(
)
(
)
x
y
x
y
I
e dxdy
dxdy
e
\pi 
\pi 
\Sigma +\Sigma +\Sigma 
\Sigma 
\Sigma 
+
\le 
+
\le 
=
-
-
=
-
-
+
-
=-
+
\int \int 
\int \int 
\int \int 
\int \int 
\int \int 

. 
\textcolor{lyred}{例}.求
(
)
xdydz
ydzdx
zdxdy
I
x
y
z
\Sigma 
+
+
=
+
+
\int \int 
,其中
(
)
:1
z
x
y
z
\Sigma 
-
=
+
\ge 
,上侧. 
解.取
(
)
1 :
z
x
y
x
y
\Sigma 
=
-
-
+
\le 
,下侧,则
I
\Sigma +\Sigma 
\Sigma 
\Sigma 
=
-
=
-
=
\int \int 
\int \int 
\int \int 

xdydz
ydzdx
zdxdy
-\Sigma 
+
+
\int \int 
,再取
(
)
2 :
z
x
y
\Sigma 
=
+
\le 
,下侧,则 
I
dv
\pi 
\pi 
-\Sigma +\Sigma 
\Sigma 
\Omega 
=
-
=
=\cdots 
=
\int \int 
\int \int 
\int \int \int 

. 
\textcolor{lyred}{例}.求
(
)
xdydz
ydzdx
zdxdy
I
x
y
z
\Sigma 
+
+
=
+
+
\int \int 
,
(
)
(
)
(
)
:
x
y
z
z
-
-
\Sigma 
+
+
=
\ge 
,上侧. 
解.取
1 : z
r
x
y
\Sigma 
=
-
-
,
(
)
(
)
:
,
x
y
z
x
y
r


-
-
\Sigma 
=
+
\ge 
+
\le 






,下侧,则 
I
xdydz
ydzdx
zdxdy
r
\pi 
\Sigma +\Sigma +\Sigma 
\Sigma 
\Sigma 
-\Sigma 
=
-
-
=
+
+
=
\int \int 
\int \int 
\int \int 
\int \int 

. 
\vspace{4pt}
{\bfseries \textcolor{lyred}{二.沿闭曲面积分为零的条件}}

\textcolor{lyred}{定理}.设G 为空间二维单连通区域,即G 内任意闭曲面所围区域均包含于G ,则 
对于G 内任意的闭曲面\Sigma ,
A dS
\Sigma 
\cdots 
=
\Leftrightarrow 
\int \int 
\rho 
\rho 

在G 内div
A =
\rho 
. 
\textcolor{lyred}{推论}.设G 为空间二维单连通区域,则在G 内,积分值
A dS
\Sigma 
\cdots 
\int \int 
\rho 
\rho 
与积分曲面无关 
\Leftrightarrow 在G 内div
A =
\rho 
处处成立. 
\textcolor{lyred}{例}.设
(
)(
)
:
2 1
z
x
y
z
\Sigma 
=
-
-
\ge 
,上侧,求I = 
(
)
yz x
y
z dydz
xz x
y
z dzdx
x y
xy x
y
z
dxdy
\Sigma 
+
+
+
+
+
+
-
+
+
\int \int 
. 
解.除原点外,
P
Q
R
x
y
z
\partial 
\partial 
\partial 
+
+
=
\partial 
\partial 
\partial 
,故取
1 :
z
x
y
\Sigma 
=
-
-
,上侧,则 
(
)
I
yzdydz
xzdzdx
x y
xy dxdy
\Sigma 
=
+
+
-
\int \int 
,又
P
Q
R
x
y
z
\partial 
\partial 
\partial 
+
+
=
\partial 
\partial 
\partial 
处处成立,再取 
(
)
2 :
z
x
y
\Sigma 
=
+
\le 
,上侧,则
(
)
x
y
I
x y
xy dxdy
x y dxdy
\pi 
\Sigma 
+
\le 
=
-
=
=
\int \int 
\int \int 
. 
\vspace{4pt}
{\bfseries \textcolor{lyred}{三.散度的物理意义}}

设A
\rho 
为一稳定不可压缩流体的流速场,则
A dS
\partial \Omega 
\Phi =
\cdots 
\int \int 
\rho 
\rho 

为单位时间内该流体由\Omega  
内部通过\partial \Omega 流向外部的流量(流出-流入),也即\Omega 内产生新流体的总量,于是, 
div
lim
div
lim
M
M
M
A
Adv
\Delta \Omega \to 
\Delta \Omega \to 
\Delta \Omega 
=
=
\Phi 
\Delta \Omega 
\Delta \Omega 
\int \int \int 
\rho 
\rho 
,代表M 处单位时间单位体积内产生 
新流体的量,故当div
M
A
>
\rho 
时,称M 为\textcolor{lyred}{源},当div
M
A
>
\rho 
时,称M 为\textcolor{lyred}{负源},或\textcolor{lyred}{汇}, 
div A
\rho 
的大小反映了源的强度,若处处div
A =
\rho 
,则称A
\rho 
为\textcolor{lyred}{无源场}. 
\textcolor{lyred}{补充练习 }
1.求
I
xzdydz
ydzdx
yzdxdy
\Sigma 
=
+
+
\int \int 
,其中
(
)
:
,
x
y
z
x z
\Sigma 
+
+
=
\ge 
,前侧. 
解.取
(
)
1 :
1,
x
y
z
z
\Sigma 
=
+
\le 
\ge 
,后侧,
(
)
2 :
1,
z
x
y
x
\Sigma 
=
+
\le 
\ge 
,下侧,则 
(
)
(
)
I
z
y dv
z
dv
\pi 
\Sigma +\Sigma +\Sigma 
\Sigma 
\Sigma 
\Sigma +\Sigma +\Sigma 
\Omega 
\Omega 
=
-
-
=
=
+
+
=
+
=
\int \int 
\int \int 
\int \int 
\int \int 
\int \int \int 
\int \int \int 


. 
2.设
()
f x 在(
)
0,+\infty 上有连续的导数,对(
)
\{
\}
, ,
:
x y z
x >
内的任意闭曲面\Sigma ,均有 
()
()
x
xf x dydz
xyf x dzdx
ze dxdy
\Sigma 
-
-
=
\int \int 
,且
()
lim
x
f x
+
\to 
=,求
()
f x . 
解.
()
()
(
)
()
()
()
div
,
,
x
x
xf x
xyf x
ze
xf
x
f x
xf x
e
'
-
-
=
+
-
-
=
,即
()
y
f x
=
满足 
x
y
y
e
x


'+
-
=




,解得
()
2x
x
e
Ce
f x
x
+
=
,由
lim
x
x
x
e
Ce
x
+
\to 
+
=,得到
C =-,故 
()
2x
x
e
e
f x
x
-
=
. 
\vspace{6pt}
{\large \bfseries \textcolor{lyred}{第11.7 节 斯托克斯公式}}

\vspace{4pt}
{\bfseries \textcolor{lyred}{一.Stokes 公式}}

\textcolor{lyred}{定义}.向量场
(
)
,
,
A
P Q R
=
\rho 
的散度为div A
A
=\nabla \cdots 
\rho 
\rho 
\rho 
,其中
,
,
x
y
z


\partial 
\partial 
\partial 
\nabla =

\partial 
\partial 
\partial 


\rho 
,又记 
rot
,
,
i
j
k
R
Q
P
R
Q
P
A
A
x
y
z
y
z
z
x
x
y
P
Q
R


\partial 
\partial 
\partial 
\partial 
\partial 
\partial 
\partial 
\partial 
\partial 
=\nabla \times 
=
=
-
-
-


\partial 
\partial 
\partial 
\partial 
\partial 
\partial 
\partial 
\partial 
\partial 


\rho 
\rho 
\rho 
\rho 
\rho 
\rho 
,称为A
\rho 
的\textcolor{lyred}{旋度}. 
\textcolor{lyred}{有向曲面边界的定向}:当观察者站在有向曲面\Sigma 选定一侧的边界上,沿着该方向 
行走时,\Sigma 位于他的左侧;或者,从\Sigma 内部看,该方向为逆时针;或者,\Sigma 的法向量 
与该方向之间满足右手法则.选择了该方向的\Sigma 的边界记为
+
\partial \Sigma . 
\textcolor{lyred}{定理}.设分段光滑有向闭曲线\Gamma 是分片光滑有向曲面\Sigma 的正向边界,若
(
)
, ,
P x y z , 
(
)
, ,
Q x y z ,
(
)
, ,
R x y z 在\Sigma 上有连续的偏导数,则 
R
Q
P
R
Q
P
Pdx
Qdy
Rdz
dydz
dzdx
dxdy
y
z
z
x
x
y
\Gamma 
\Sigma 




\partial 
\partial 
\partial 
\partial 
\partial 
\partial 


+
+
=
-
+
-
+
-






\partial 
\partial 
\partial 
\partial 
\partial 
\partial 






\int 
\int \int 
\nabla 
,即 
(
)
rot
rot
n
A dr
A dS
A e
dS
\Gamma 
\Sigma 
\Sigma 
\cdots 
=
\cdots 
=
\cdots 
\int 
\int \int 
\int \int 
\rho 
\rho 
\rho 
\rho 
\rho 
\rho 
\nabla 
,其中
ne\rho 为\Sigma 的单位正法向量. 
\textcolor{lyred}{注}.
A dr
\Gamma 
\cdots 
\int 
\rho 
\rho 
\nabla 
称为向量场A
\rho 
沿曲线\Gamma 的\textcolor{lyred}{环流量}. 
\textcolor{lyred}{例}.证明:(1)
(
)
div rot
A =
\rho 
;(2)
(
)
rot grad
f
=
\rho 
. 
证.(1)对于空间任意具有光滑边界的有界闭区域\Omega ,均有 
(
)
div rot
rot
A dv
A dS
A dr
\Omega 
\partial \Omega 
\partial \partial \Omega 
=
\cdots 
=
\cdots 
=
\int \int \int 
\int \int 
\int 
\rho 
\rho 
\rho 
\rho 
\rho 

\nabla 
,故
(
)
div rot
A =
\rho 
; 
(2)对于空间任意具有光滑边界的光滑有向曲面\Sigma ,均有 
(
)
rot grad
grad
f
dS
f dr
df
\Sigma 
\partial \Sigma 
\partial \Sigma 
\cdots 
=
\cdots 
=
=
\int \int 
\int 
\int 
\rho 
\rho 
\nabla 
\nabla 
,故
(
)
rot grad
f
=
\rho 
,证毕. 
\textcolor{lyred}{例}.求I
zdx
xdy
ydz
\Gamma 
=
+
+
\int \nabla 
,其中\Gamma 为
x
y
z
+
+
=被三个坐标面截下的三角形的 
边界,取向与三角形的上侧法向之间符合右手法则. 
解.
(
)
(
)
(
)
rot
, ,
1,1,1
1,1,1
n
I
z x y
e dS
dS
dS
\Sigma 
\Sigma 
\Sigma 
=
\cdots 
=
\cdots 
=
=
\cdots 
=
\int \int 
\int \int 
\int \int 
\rho 
; 
或者,
(
)
rot
, ,
xy
D
I
z x y
dS
dydz
dzdx
dxdy
dxdy
d\sigma 
\Sigma 
\Sigma 
\Sigma 
=
\cdots 
=
+
+
=
=
=
\int \int 
\int \int 
\int \int 
\int \int 
\rho 
. 
\textcolor{lyred}{例}.求I
ydx
zdy
xdz
\Gamma 
=
+
+
\int \nabla 
,其中
(
)
:
x
y
z
x
y
x
y

+
+
=
+
\Gamma +=

,从x 轴正向看为顺时针. 
解.
(
)
(
)(
)
1, 1,0
rot
, ,
1, 1, 1
2 2
n
I
y z x
e dS
dS
dS
\pi 
\Sigma 
\Sigma 
\Sigma 
--
=
\cdots 
=
---
\cdots 
=
=
\int \int 
\int \int 
\int \int 
\rho 
. 
\textcolor{lyred}{例}.求
(
)
(
)
(
)
I
y
z
dx
z
x
dy
x
y
dz
\Gamma 
=
-
+
-
+
-
\int \nabla 
,其中
:
x
y
x
y
z

+
=
\Gamma +
+
=

,从z 轴的 
正向看为逆时针. 
解.
(
)
(
)
(
)
rot
,
,
1,1,1
n
I
A e dS
y
z z
x x
y
dS
x
y
z dS
\Sigma 
\Sigma 
\Sigma 
=
\cdots 
=
-
+
+
+
\cdots 
=-
+
+
=
\int \int 
\int \int 
\int \int 
\rho \rho 
2 1
3 x
y
dS
dxdy
\pi 
\Sigma 
+
\le 
-
=-
=-
\int \int 
\int \int 
. 
\textcolor{lyred}{例}.求I
xyzdz
\Gamma 
=\int \nabla 
,其中
:
y
z
x
y
z
-
=

\Gamma 
+
+
=

,从z 轴正向看为逆时针. 
解.
(
)
(
)
(
)
rot 0,0,
,
,0
0, 1,1
n
I
xyz
e dS
xz
yz
dS
yzdS
\Sigma 
\Sigma 
\Sigma 
=
\cdots 
=
-
\cdots 
-
=
=
\int \int 
\int \int 
\int \int 
\rho 
x
y
x
y
x
y
y
dxdy
y dxdy
y
dx dy
'
'
+
\le 
+
\le 
+
\le 

'
'
'
\cdots 
=
=
\cdots 
=




\int \int 
\int \int 
\int \int 
(
)
x
y
x
y
dx dy
\pi 
'
'
+
\le 
'
'
'
'
+
=
\int \int 
. 
\textcolor{lyred}{例}.求
I
xzdx
x dy
\Gamma 
=
+
\int \nabla 
,其中
:
z
x
y
x
y
ax
=
+
\Gamma 
+
=

,从z 轴正向看为顺时针. 
解.
(
)
(
)
(
)
2 ,2 , 1
rot
,
,0
0, ,2
n
x
y
I
xz x
e dS
x
x
dS
x
y
\Sigma 
\Sigma 
-
=
\cdots 
=
\cdots 
=
+
+
\int \int 
\int \int 
\rho 
(
)
(
)
xy
xy
xy
D
D
D
xy
x
x
y dxdy
xy
x dxdy
xdxdy
x
y
-
+
+
=
-
=-
=
+
+
\int \int 
\int \int 
\int \int 
(
)
xy
xy
D
D
x
a
a dxdy
a
dxdy
a
\pi 
-
-
+
=-
=-
\int \int 
\int \int 
. 
\textcolor{lyred}{例}.求
I
y dx
z dy
x dz
\Gamma 
=
+
+
\int \nabla 
,其中
: z
a
x
y
x
y
ax
=
-
-
\Gamma 
+
=

,从z 轴正向看为逆时针. 
解.
(
)
(
)
rot
,
,
2 , 2 , 2
,
,
n
x y z
I
y
z
x
e dS
z
x
y
dS
a a a
\Sigma 
\Sigma 


=
\cdots 
=
-
-
-
\cdots 
=




\int \int 
\int \int 
\rho 
(
)
xy
D
adxdy
xz
xy
yz dS
xzdS
x a
x
y
a
a
a
a
x
y
\Sigma 
\Sigma 
-
+
+
=-
=-
-
-
\cdots 
=
-
-
\int \int 
\int \int 
\int \int 
xy
xy
xy
D
D
D
a
a
a
xdxdy
x
dxdy
dxdy
a
\pi 


-
=-
-
+
=-
\cdots 
=-




\int \int 
\int \int 
\int \int 
. 
\vspace{4pt}
{\bfseries \textcolor{lyred}{二.曲线积分与路径无关的条件}}

\textcolor{lyred}{定义}.称空间区域G 为\textcolor{lyred}{一维单连通区域},若G 内每条闭曲线C 均是包含于G 内的 
某个曲面的边界. 
\textcolor{lyred}{定理}.设G 为空间一维单连通区域,
(
)
, ,
P x y z ,
(
)
, ,
Q x y z ,
(
)
, ,
R x y z 在G 内均有 
连续的偏导数,则下列四个条件等价: 
(1)在G 内曲线积分Pdx
Qdy
Rdz
\Gamma 
+
+
\int 
与路径无关; 
(2)对于G 内任意闭曲线C ,均有
C
Pdx
Qdy
Rdz
+
+
=
\int \nabla 
; 
(3)在G 内
(
)
rot
,
,
P Q R =
\rho 
;(4) 在G 内微分形式Pdx
Qdy
Rdz
+
+
可积. 
\textcolor{lyred}{注}.此时, (
)
(
)
(
)
, ,
,
,
, ,
x y z
x
y
z
u x y z
Pdx
Qdy
Rdz
=
+
+
\int 
是它在G 内的一个原函数. 
\textcolor{lyred}{例}.求
(
)
(
)
(
)
I
x
yz dx
y
xz dy
z
xy dz
\Gamma 
=
-
+
-
+
-
\int 
,其中\Gamma 为
cos
sin
x
y
z
\varphi 
\varphi 
\varphi 
=


=

=

上从 
(
)
1,0,0
A
到
(
)
1,0,2
B
\pi 的一段有向弧. 
解.
(
)
rot
,
,
x
yz y
xz z
xy
-
-
-
=
,故
\approx 
(
)
1 0
AB
I
z
dz
\pi 
\pi 
=
=
-\cdots 
=
\int 
\int 
. 
\vspace{4pt}
{\bfseries \textcolor{lyred}{三.曲线积分的基本定理}}

\textcolor{lyred}{定理}.设在G 内Pdx
Qdy
Rdz
du
+
+
=
,则对G 内分段光滑有向弧段
\approx 
L
AB
=
,均有 
\approx 
\approx 
()
()
AB
AB
Pdx
Qdy
Rdz
du
u B
u A
+
+
=
=
-
\int 
\int 
. 
\textcolor{lyred}{向量形式}.若
grad
F
u
=
\rho 
,则
\approx 
\approx 
()
()
AB
AB
F dr
du
u B
u A
\cdots 
=
=
-
\int 
\int 
\rho 
\rho 
. 
\textcolor{lyred}{例}.设\Gamma 为光滑闭曲线,则
(
)
grad sin
x
y
z
dr
\Gamma 
+
+
\cdots 
=




\int 
\rho 
\nabla 
. 
\textcolor{lyred}{例}.设(
)
0,0,1
A
,
(
)
1,2,2
B
,
r
x
y
z
=
+
+
,求
(
)
\approx 
AB
I
r
xdx
ydy
zdz
=
+
+
\int 
. 
解.
\approx 
\approx 
\approx 
(
)
AB
AB
AB
I
r dr
r dr
dr
=
=
=
=
-
=
\int 
\int 
\int 
. 